%Plantilla basada en "Template for Masters / Doctoral Thesis" (plantilla disponible en writeLaTex) que subió LaTeXTemplates.com

\documentclass[11pt]{book}
\usepackage[paperwidth=17cm, paperheight=22.5cm, bottom=2.5cm, right=2.5cm]{geometry}
\usepackage[utf8]{inputenc}
\usepackage{amssymb,amsmath,amsthm} %paquete para símbolo matemáticos
\usepackage{amsmath,amsfonts}
\usepackage{bbm}
\usepackage{physics}
\usepackage[spanish]{babel}
\usepackage[utf8]{inputenc} %Paquete para escribir acentos y otros símbolos directamente
\usepackage{enumerate}
\usepackage{graphicx}
\usepackage{xr}
\usepackage{verbatim}
%\usepackage{subfig} %para poner subfiguras
\graphicspath{{Img/}} %En qué carpeta están las imágenes
\usepackage[nottoc]{tocbibind}
\usepackage[pdftex,
            pdfauthor={Germán Eduardo Osorio Leiva},
            pdftitle={Revisión de los fundamentos de la mecánica estadística desde la perspectiva de la teoría cuántica de la información},
            pdfsubject={Física},
            pdfkeywords={entrelazamiento, mecánica estadística},
            pdfproducer={Latex con hyperref},
            pdfcreator={pdflatex}]{hyperref}
            
\newtheorem{theorem}{Teorema}[section]
\newtheorem{corollary}{Corolario}[theorem]
\newtheorem{lemma}[theorem]{Lema}
\theoremstyle{definition}
\newtheorem{definition}{Definición}[section]




\begin{document}

%----------------------------------------------------------------------------------------
%	COMANDOS PERSONALIZADOS
%----------------------------------------------------------------------------------------

%SI TU TESIS TIENE TEOREMAS Y DEMOSTRACIONES, PUEDES DESCOMENTAR Y USAR LOS SIGUIENTES COMANDOS

%\proofname}{Demostración}
%\providecommand{\norm}[1]{\lVert#1\rVert} %Provee el comando para producir una norma.
%\providecommand{\innp}[1]{\langle#1\rangle} 
%\newcommand{\seno}{\mathrm{sen}}
%\newcommand{\diff}{\mathrm{d}}

%\newtheorem{teo}{Teorema}[section] 
%\newtheorem{cor}[teo]{Corolario}
%\newtheorem{lem}[teo]{Lema}

%\theoremstyle{definition}
%\newtheorem{dfn}[teo]{Definición}

%\theoremstyle{remark}
%\newtheorem{obs}[teo]{Observación}

%\allowdisplaybreaks


%----------------------------------------------------------------------------------------
%	PORTADA
%----------------------------------------------------------------------------------------

\title{Revisión de los fundamentos de la mecánica estadística desde la perspectiva de la teoría cuántica de la información} %Con este nombre se guardará el proyecto en writeLaTex

\begin{titlepage}
\begin{center}

\textsc{\Large Universidad de los Andes}\\[4em]

%Figura
\begin{figure}[h]
\begin{center}
\includegraphics[scale=0.045]{ellogo.png}
\end{center}
\end{figure}

\vspace{4em}

\textsc{\huge \textbf{Revisión de los fundamentos de la mecánica estadística desde la perspectiva de la teoría cuántica de la información}}\\[4em]

\textsc{\large Tesis}\\[1em]

\textsc{para obtener el título de}\\[1em]

\textsc{FÍSICO}\\[1em]

\textsc{presenta}\\[1em]

\textsc{\Large 	GERMÁN EDUARDO OSORIO LEIVA}\\[1em]
		
\textsc{\large Asesor: Alonso Botero Mejía}

\end{center}

\vspace*{\fill}
\textsc{Bogotá, D.C. \hspace*{\fill} 2017}

\end{titlepage}


%----------------------------------------------------------------------------------------
%	DECLARACIÓN
%----------------------------------------------------------------------------------------

\thispagestyle{empty}
\vspace*{8em}
\textit{"Y así pensé yo que las ciencias de los libros, por lo menos aquellas cuyas razones son sólo probables y carecen de demostraciones, habiéndose compuesto y aumentado poco a poco con las opiniones de varias personas diferentes, no son tan próximas a la verdad como los simples razonamientos que un hombre de buen sentido puede hacer... "}
\begin{flushright}
René Descartes
\end{flushright}


%\vspace*{\fill}
%\begingroup

%\centering

%\hspace{3em}

%\textsc{AUTOR}

%\vspace{5em}

%\rule[1em]{20em}{0.5pt} % Línea para la fecha

%\textsc{Fecha}
 
%\vspace{8em}

%\rule[1em]{20em}{0.5pt} % Línea para la firma

%\textsc{Firma}

%\endgroup
%\vspace*{\fill}


%----------------------------------------------------------------------------------------
%	DEDICATORIA
%----------------------------------------------------------------------------------------

\pagestyle{empty}
\frontmatter
RESUMEN

%----------------------------------------------------------------------------------------
%	AGRADECIMIENTOS
%----------------------------------------------------------------------------------------


\chapter*{Agradecimientos}
Agradezco a mi padre que me ha apoyado en todos mis proyecto confiando en mí y dándome la fortaleza necesaria para las situaciones complicadas. También, siento mucha gratitud hacía mi madre que con su energía y alegría ha hecho de la vida un paseo muy colorido. De la misma manera, hay muchas personas que saben lo importante que son para mí y han hecho de este viaje muy ameno. Estoy muy agradecido con mi asesor Alonso Botero por su orientación.
%\markboth{INTRODUCCIÓN}{AGRADECIMIENTOS} % encabezado 




%----------------------------------------------------------------------------------------
%	PREFACIO
%----------------------------------------------------------------------------------------

%\chapter*{Prefacio}

%\pagestyle{plain}
%\markboth{PREFACIO23}{PREFACIO} % encabezado 

%PUEDEN QUITAR ESTA PARTE


%----------------------------------------------------------------------------------------
%	TABLA DE CONTENIDOS
%---------------------------------------------------------------------------------------

\tableofcontents


%----------------------------------------------------------------------------------------
%	TESIS
%----------------------------------------------------------------------------------------
\mainmatter %empieza la numeración de las páginas
\pagestyle{headings}

%  Incluye los capítulos en el folder de capítulos

\chapter{Introducción} 
%El escrito que se presentará a continuación, quiere centrarse en exponer cómo los fundamentos de la mecánica estadística tiene problemas y su posible resolución gracias al trabajo de Popescu et al \cite{Popescu2006}. Para lograr esto se espera analizar los fundamentos que actualmente tiene la mecánica estadística, mirar las perspectivas más importantes siguiendo una exposición histórica. Para luego poder exponer el enfoque dado por Popescu et al se reproducirán los cálculos del artículos. Con el objetivo de exponer claramente las ideas aportadas en el artículos.

La mecánica estadística muestra gran poder para comprender los fenómenos causados por sistemas de muchas partículas. Las predicciones teóricas concuerdan con los experimentos, pero el marco teórico en el que se fundamenta tiene problemas conceptuales que dejan a la teoría en un cimiento débil.  Para exhibir las incongruencias internas de los fundamentos. Se exploran las ideas de Boltzmann y de Gibbs, porque sus enfoques tienen un gran impacto en la creación del marco conceptual de la mecánica estadística.
\\ 
En la época de Botlzmann se encontraba la termodinámica como una teoría establecida. Aunque, no se tenían una óptica donde se le diera significado microscópico a la termodinámica. Boltzmann buscando una conexión entre la termodinámica y  el mundo microscópico plantea el teorema H. Este teorema muestra que al pasar el tiempo la función H decrece. La relación con la entropía termodinámica se da al multiplicar la función H por una constante -K, entonces se tiene que la entropía aumenta con el tiempo. El problema con la formulación hecha por Boltzmann, según Jaynes, es que la función H solo entrega la entropía termodinámica  cuando no hay interacción entre partículas.
\\
Unos años después, Gibbs postuló su noción de ensambles. Jaynes muesta que la entropía de Gibbs, para todos los casos, es la entropía termodinámica. Aunque sus resultados concuerden con la termodinámica, la entropía de Gibbs no puede dar explicación a la flecha del tiempo, porque esta entropía se conserva a diferencia de la de Boltzmann. Pese a esto el enfoque de Gibbs se muestra poderoso para resolver muchos problemas, pero Gibbs no explicó ¿cómo es posible hablar de un ensamble de posbiles estados microscópicos, estando en un solo estado microscópico? En la literatura se resuelve esta pregunta usando la hipótesis ergódica. Esto acarrea problemas porque no todos los sistemas son ergódicos; y el tiempo para que sistemas ergódicos pasen por todo el espacio son extremadamente grandes. 
\\
Partiendo de los planteamientos realizados por Popescu, Short y Winter, en el artículo “Entanglement and the foundations of statistical mechanics” de \cite{Popescu2006}. Donde se estudia la idea de tipicidad canónica con herramientas de la teoría de la información cuántica. Argumentando que cuando se tiene un sistema cuántico aislado compuesto por el subsistema y el ambiente, siendo el ambiente mucho más grande que el subsistema, se asegura que, aunque se tenga un estado puro de todo el sistema cuántico el subsistema se comportará como si el sistema estuviese en un estado máximamente mezclado. El autor formaliza esta idea dando unas cotas sobre la probabilidad de que haya un subsistema que no se comporte como un estado canónico. Estas cotas decaen de manera exponencial con respecto a la dimensión del sistema. Entonces se postula un nuevo principio, el principio de probabilidades iguales aparentes; por lo tanto, el subsistema tiene la misma probabilidad de estar en cualquier estado admisible. 
\\
Por otra parte, los planteamientos; igualmente presentados por los mismos autores junto con  Linden en el artículo “Quantum mechanical evolution towards thermal equilibrium” \cite{LindenPaper}  donde se expone en base a la tipicidad canónica cómo los sistemas llegan al equilibrio; partiendo de la explicación sobre qué significa que un sistema llegue al equilibrio, manifestando  que un sistema se dice estar en equilibrio cuando pasa mucho tiempo cerca de un estado particular. También se necesita una condición de independencia sobre el estado inicial del ambiente y el estado inicial del subsistema, porque el estado de equilibrio no debe depender de las condiciones iniciales que tenga el sistema.  La formulación matemática exhibida en este artículo muestra que los estados genéricos y los estados fuera del equilibrio, del subsistema, llegan a termalizarse a un estado específico, siendo este estado el promedio temporal del estado del subsistema.
\\
Desde este horizonte de sentido, en el presente trabajo, se realizará una revisión crítica del estado del arte de las ideas modernas que han dado lugar a una reformulación de los fundamentos estadísticos de la termodinámica a partir de las herramientas que proporciona la teoría de la información cuántica, así mismo, se consideran algunos ejemplos de aplicaciones de estas formulaciones a sistemas de lejos del equilibrio.
\\
En este sentido, se trata de vislumbrar las ideas sobre la esencia de la formulación estadística según la tipicidad canónica de estados cuánticos de espacios  de Hilbert  de altas dimensiones replicando los cálculos y desglosando los pasos conceptuales con importancia física , concentrándose en las ideas de Popescu et al;  y  revisar detenidamente los cálculos e ideas nuevas  sobre las aplicaciones de la termodinámica en sistemas fuera de equilibrio dadas por la teoría de la información cuántica, comparando con los avances teóricos  hechos en años anteriores  sin estas herramientas,
\\
Es así como de manera puntual, en el primer capítulo se esboza primero, algunas orientaciones generales haciendo un breve recuento histórico desde la termodinámica, la teoría cinética; luego los fundamentos de la mecánica estadística desde la perspectiva de Boltzmann y Gibbs; después  el ensamble canónico, el teorema de Liouville  y sus consecuencias, la ergocidad  y finalmente ,el principio de máxima entropía de Jaynes     
\\
En el segundo capítulo, se plantea en una primera instancia el entrelazamiento y la mecánica estadística, posteriormente las herramientas de la teoría cuántica, luego el operador de densidad y operador de densidad reducida, después, distancia de traza, algunas definiciones, la idea conceptual, la formulación matemática, el lema de Levy, la demostración del principio general canónico y por último el significado físico.
\\
En el tercer capítulo se presenta primero la evolución hacia el equilibrio, luego la especificación del equilibrio para terminar con la equilibración de los sistemas.
\\
Finalmente se presentan algunas conclusiones, producto del desarrollo del presente trabajo. 
%\\

%Nuevas propuestas han salido para resolver estos conflictos, como la perspectiva de Popescu et al. Este nuevo enfoque permite conciliar los puntos de vista anteriores usando herramientas de la teoría cuántica de la información. Esta formulación muestra que, si se tiene un sistema, compuesto por el ambiente y el subsistema, localmente (en el subsistema) se puede describir con una densidad canónica. Entonces hace válida la densidad canónica sin hacer uso de los ensambles, además las probabilidades no se introducen de una manera ajena a la teoría física. Se demuestra que el subsistema está en un estado equiprobable sin necesidad de mostrar ergodicidad en el sistema. También muestra de forma implícita la maximización de entropía por el entrelazamiento. Así es como esta perspectiva muestra una imagen más clara y fuerte del marco teórico. Linden et al \cite{LindenPaper} mostraron que se puede explicar la termalización de los sistemas con ayuda de esta perspectiva, lo cual hace a la nueva formulación útil e interesante para futuras investigaciones.
%\\
%El orden que seguirá el siguiente escrito seguirá el esqueleto dado en esta introducción. Primero se mostrará el estado del arte concerniente a los fundamentos de la mecánica estadística. Luego se pasará a exhibir el enfoque de Popescu et al y finalizando con  la construcción basada en esta perspectiva. 


\thispagestyle{empty}
\chapter{El entrelazamiento y la mecánica estadística}

La idea principal que se quiere mostrar es cómo se puede reemplazar el postulado de probabilidades iguales por un principio canónico general basado en el entrelazamiento cuántico, basado en el entrelazamiento del sistema y el ambiente.

La mecánica clásica nos habla de un sistema físico definido que para todos los tiempos está especificado. Este sistema evolucionan de manera determinista. Pero lo que sorprende al tratar sistemas termodinámicos es que aunque se hable de un sistema clásico este puede mostrar propiedades que dependan de promedios estadísticos. La conexión que hay entre el determinismo y estas probabilidades es una discusión que lleva desde los inicios de la termodinámica. Aunque el planteamiento que se mostrará no necesita de los métodos típicos de la mecánica estadística como por ejemplo: aleatoriedad subjetiva, promedio sobre el ensamble o promedio temporal; Esto le da una fuerza a estas ideas ya que no debe entrar en problemas de ergodicidad. 

Este nuevo enfoque se tiene un universo que está compuesto por el sistema y el ambiente. Dando como condición que el ambiente sea los suficientemente grande. Este universo está descrito por un estado cuántico puro ( se conoce el estado de manera exacta) que obedece una restricción global, que el sistema alcance el equilibrio térmico por medio de la interacción mutua (termalización) es un producto del entrelazamiento del sistema y el ambiente (Popescu et al,2006). Más adelante se le dará una definición más rigurosa que ayudará a dar cotas para la expresión "ambiente suficientemente grande". Con este enfoque se quiere formular un principio canónico general  el sistema estará termalizado para casi todos los estados puros del universo, dando límites cuantitativos. La restricción que se impone no es una específica esto generaliza los resultados tradicionales dados en la literatura donde se toma por restricción la energía (cualquier libro, año).
Estos resultados no miran la evolución del sistema pero debido a que la mayoría de los estados del universo están termalizados se preve que no importa el estado inicial cualquier evolución llevará al estado a uno en el equilibrio.

(
Lo que mi asesor no me ha respondido.
)

Ya poniendo las ideas más explícitas se supone tener un sistema cuántio aislado y grande que se llamará el universo, que se descompone en dos partes el sistema $S$ y el ambiente $E$. La dimensión del ambiente es mucho maś grande que la del sistema $S$. Además se le impone una restricción global al universo llamada $R$. Desde la mecánica cuántica estas condiciones se pueden poner como restricciones en el espacio de Hilbert, restricción de los estados posibles:

\begin{equation}
\mathcal{H_{R}}\subseteq \mathcal{H_{S}}\otimes \mathcal{H_{E}},
\end{equation}

Donde $\mathcal{H_{S}}$ y $\mathcal{H_{E}}$ son los espacios de Hilbert del sistema y el ambiente con dimension $d_{\mathcal{S}}$  y $d_{\mathcal{E}}$ respectivamente. Es bueno recalcar que $R$ es una restricción arbitraria generalmente se toma como la energía del universo.
Ahora se define el estado equiprobable del universo bajo $R$ como

\begin{equation}
\mathcal{E_{R}} = \frac{1}{d_{R}} \mathbbm{1_{\mathcal{R}}},
\end{equation}

Donde $\mathbbm{1_{\mathcal{R}}}$ es el operador identidad (proyección) sobre el espacio de Hilbert  $\mathcal{H_{R}}$  que tiene dimensión $d_{\mathcal{R}}$. Esto se relaciona con el principio de probabilidades iguales porque este es el estado máximamente mezclado en $\mathcal{H_{R}}$ por ser así todos los estados bajo la restricción de R tienen la misma probabilidad de salir.

Definimos $\Omega_{\mathcal{S}}$ como el estado canónico que está restringido por $R$ cuando el universo se encuentra en el estado $\mathcal{E_{R}}$.Esto significa que si se hace una traza parcial al universo sobre el ambiente da como resultado el estado canónico:
\begin{equation}
 \Omega_{\mathcal{S}} =Tr_{\mathcal{E}} \mathcal{E_{R}}.
\end{equation}

Aquí se hace un supuesto importante y es que el universo está en un estado puro $\ket{\phi}$ y no en un estado mixto $\mathcal{E_{R}}$, esto quiere decir que se conoce todo lo que es permitido por la mecánica cuántica del universo si estuviese en un estado mixto significaría que nosotros no tenemos toda la información que se pudiese tener. Ahora lo que se quiere ver es que aun que el estado del universo sea puro el estado reducido del sistema,

\begin{equation}
\rho_{S}=Tr_{E}\ket{\phi}\bra{\phi} 
\end{equation}

se acerca al estado canónico para la gran mayoría de los casos 

\begin{equation}
\rho_{S} \approx \Omega_{S}.
\end{equation}

O sea que para casi todos los estados puros $\ket{\phi} \in \mathcal{H_{R}}$ del universo el sistema se comporta como si el universo estuviese en el estado mixto equiprobable $\mathcal{E_{R}}$. Siendo este el principio general canónico.

Clarificando la idea, el estado canónico $\Omega_{S}$ del sistema es el estado del sistema cuando el universo se encuentra en el estado equiprobable $\mathcal{E_{R}}$. Se puede interpretar el principio general canónico como un principio que estipula que las probabilidades iguales del sistema son aparentes. Porque para casi cualquier estado del universo que sea puro, un subsistema de este universo que cumpla con ser lo suficientemente pequeño se comporta como si el universo estuviese en el estado equiprobable $\mathcal{E_{R}}$. Cabe recordar que aún no se ha especificado la restricción $R$ entonces todo este análisis es general, la restricción no necesariamente debe ser la energía u otras cantidades que se conserven. Esto hace que $\Omega_{S}$ no  deba ser obligatoriamente es el estado canónico usual, puede ser el gran canónico o cualquier otro que sea acorde con la restricción. Este principio puede ser de utilidad cuando la interacción entre el ambiente y el sistema no es débil como siempre se toma.

\section{formulación matemática}

Ahora se procederá a especificar lo dicho anteriormente en un contexto matemático. Lo primordial es decir cuál será la distancia que usaremos para darle un sentido de cercanía a los estados $\rho_{S}$ y $\Omega_{S}$. La distancia a usar es una bastante conocida en la teoría cuántica de la información, la distancia de traza. Esta se define como:
\begin{equation}
D(\rho_{S}, \Omega_{S})= \frac{1}{2} Tr |\rho_{S} -\Omega_{S}|=\frac{1}{2}Tr\sqrt{(\rho_{S} -\Omega_{S})^{\dag}(\rho_{S} -\Omega_{S})}.
\end{equation}

Esta medida cuantifica que tan difícil es diferenciar $\rho_{S}$ y $\Omega_{S}$ por mediciones cuánticas. Un poco de notación $\langle . \rangle$ es el promedio sobre todos los estados $\ket{\phi} \in \mathcal{H_{R}}$ de acuerdo a la medida estandar(unitariamente invariante). Esta medida se usa para hallar volúmenes de conjuntos de estados. Luego por esto se sabe que $\Omega_{S} = \langle \rho_{S} \rangle $.

Entonces el teorema central de este capítulo es:

\textbf{Teorema:}Para un estado escogido de manera aleatoria $\ket{\phi} \in \mathcal{H_{R}} \subseteq \mathcal{H_{S}} \otimes \mathcal{H_{E}} $ y un $\epsilon > 0$ arbitrario, la distancia entre la matriz densidad reducida del sistema $\rho_{S}=Tr(\ket{\phi} \bra{\phi})$  y el estado canónico $\Omega_{S}=Tr \mathcal{E_{R}}$ esta dado probabilísticamente por 

\begin{equation}
Prob[  \norm{\rho_{S} -\Omega_{S}}_{1} \geq \eta ] \leq \eta',
\end{equation}

Donde 

\begin{equation}
\norm{\rho_{S}-\Omega_{S}}_{1}=2 D(\rho_{S}, \Omega_{S})
\end{equation}

\begin{equation}
\eta= \epsilon + \sqrt{ \frac{d_{S}}{d_{E}^{eff}} },
\end{equation}

\begin{equation}
\eta'=2exp(-Cd_{R}\epsilon^{2}).
\end{equation}

y las constantes son: $ C=(18\pi^{3})^{-1}, d_{R} = dim \mathcal{H_{R}}, d_{S} = dim \mathcal{H_{S}} $. $d_{E}^{eff}$ es la medida efectiva del tamaño del ambiente,
\begin{equation}
d_{E}^{eff}= \frac{1}{Tr \Omega_{E}^{2}} \ge \frac{d_{R}}{d_{S}}.
\end{equation}

Donde $\Omega_{E}= Tr_{S} \mathcal{E_{R}}$. Ambas cantidades $\eta $ y $\eta'$ serán pequeñas esto implica que el estado estará cercano al estado canónico con alta probabilidad cuando la dimensión efectiva del ambiente sea mucho más grande que la del sistema ( es decir $d_{E}^{eff} >> d_{S}$) y  $d_{R}\epsilon^2>>1>>\epsilon$. Esta última condición se puede asegurar cuando el espacio acesible total es grande ( es decir $d_{R}>>1$), escogiendo $\epsilon=d_{r}^{-\frac{1}{3}}$.\\

Ya con esto se le da un significado cuantitativo de lo que se ha ido explicando hasta ahora. Además también se dan unas cotas explícitas a cuando se habla de la "mayoría" de los estados. Se tiene una cota exponencialmente pequeña  sobre la probabilidad de encontrar un estado lejano del canónico. (Dar comentario sobre lo intuitivo de las cotas)
\section{Lema de Levy}
El teorema anterior tiene como base principal el lema de Levy (referencias). Este dice que al seleccionar un punto $\phi$ aleatoriamente de una hiperesfera de dimensión alta y que $f(\phi)$ no cambie muy rápido, entonces $f(\phi) \approx \langle f \rangle $ con alta probabilidad, Más exactamente:

\textbf{Lema de Levy:} Dada una función $f: \mathbb{S}^d \to \mathbb{R} $ definida en la hiperesfera $\mathbb{S}^d$, y un punto $\phi \in \mathbb{S}^d $ es escogido de manera aleatoria uniforme, 
\begin{equation}
Prob[|f(\phi)- \langle f \rangle| \ge \epsilon ] \le 2exp(-\frac{2C(d+1)\epsilon^2}{\eta^2})
\end{equation}
donde $\eta$ es la constante de Lipschitz de $f$, dado por $\eta= \sup|\nabla f|$ y $C=(18\pi^3)^-1$.







\thispagestyle{empty}
\chapter{El Entrelazamiento y La Mecánica stadística} \label{cap:2}



La mecánica clásica nos habla de un sistema físico definido que para todos los tiempos está especificado. Este sistema evolucionan de manera determinista dadas las ecuaciones de movimiento. Pero lo que sorprende al tratar con  sistemas termodinámicos es que aunque se hable de un sistema clásico este puede mostrar propiedades que dependan de promedios estadísticos \cite{CallenThermo}. La conexión que hay entre el determinismo y estas probabilidades es una discusión que no se ha podido solucionar dado a que han existido varias soluciones pero todas con sus fallas . Los métodos típicos requieren hablar de promedios de ensamble, promedios temporales, aleatoriedad. Todos estos conceptos siguen siendo muy debatidos y parecería que no solucionará ninguna pregunta seguir con ellos \cite{TodaStat}. Por eso se ha estado buscando nuevos fundamentos para la mecánica estadística.
\\
La mecánica estadística tiene como base el postulado de probabilidades iguales; este viene dado por la ignorancia subjetiva que el observador del sistema tiene. Aunque la mecánica estadística no tiene problemas al comparar sus resultados teóricos con los experimentales, el cimiento filosófico en el que reposa da mucho de que hablar dando así posturas diferentes que desde los inicios de su teoría no han podido ser unificados. En este capítulo se expondrá cómo Popescu et al. en \cite{Popescu2006} muestran una posible luz sobre el problema de unificar las perspectivas de la mecánica estadística. La idea principal que se quiere mostrar es cómo se puede reemplazar el postulado de probabilidades iguales por un principio canónico general basado en el entrelazamiento cuántico, al poner el entrelazamiento cuántico como nuevo cimiento ya no se tiene probabilidades subjetivas sino objetivas dadas por la teoría  cuántica. Este nuevo enfoque permite evitar problemas con la ignorancia subjetiva pero también esquiva el problema de ergodicidad, gracias a esto se puede ver un fundamento más claro y sólido. 
\\
Este capítulo solo se enfocará en mostrar cómo la mayoría de los estados del universo están termalizados sin hablar en ningún momento de la evolución de estos estados. Si en general los estados del universo están termalizados se esperaría que cualquier evolución lleve los estados al equilibrio pero en el siguiente capítulo se darán especificaciones de esta conjetura junto con detalles que hacen ver las complicaciones de formalizar las ideas intuitivas que se tienen sobre el equilibrio.
\\

\subsection{Definiciones}

Antes de seguir sería preciso dar la notación que se usará a lo largo de este capítulo y el siguiente. Cuando se tenga un sistema cuántico grande descrito por un espacio de Hilbert $\mathcal{H}$  este se llamará el universo. Este universo será dividido en  dos subsistemas. El primero se llamará el sistema $S$  y el segundo se le dará el nombre de ambiente $E$, en ocasiones $S$ se le dirá subsistema y $E$ se le podrá llamar baño estos nombres son equivalentes a los anteriores. Esto lleva a descomponer $\mathcal{H}$ como $\mathcal{H}=\mathcal{H}_{S} \otimes \mathcal{H}_{E}$, con dimensiones $d_{S}$ y $d_{E}$ respectivamente, se supone que la dimensión del ambiente es mucho mayor que la del sistema. Para evitar espacios de dimensión infinita se introduce un tope para altas energías y así mantener la dimensión finita. Además se eliminarían términos de interacción del Hamiltoniano que lo lleven a subespacios no permitidos. Aún no se ha especificado nada sobre el subsistema o el ambiente (Excepto las proporciones de sus dimensiones)  esto permite decir que cualquier descomposición del espacio de Hilbert especifica un subsistema y un baño. El subsistema S puede ser cualquier cosa desde una partícula hasta el conjunto de partículas distribuidas por todo el baño.
\\
El estado global puro del universo se escribirá como $\ket{\phi(t)}$ en un tiempo t y su matriz de densidad se escribirá $\rho(t)=\ket{\phi(t)}\bra{\phi(t)}$. El estado del sistema se encontrará al hacer una traza parcial del ambiente sobre el estado del universo $\rho_{S}= \Tr _{B} \rho(t)$; similarmente el estado del ambiente está dado por $\rho_{B}=\Tr _{S} \rho(t)$ \cite{WildeInformation}. El promedio temporal del universo es:
\begin{equation}
\omega= \langle \rho(t) \rangle= \lim_{\tau \to \infty} \frac{1}{\tau} \int_{0}^{\tau} \rho (t)dt
\end{equation}
de manera análoga $\omega_{S}$ y $\omega_{B}$ son el promedio temporal para el sistema y el ambiente respectivamente \cite{TodaStat}. También se tiene el promedio $\langle . \rangle_{\phi}$ que es sobre todos los estados $\ket{\phi} \in \mathcal{H}_{R}$ de acuerdo a la medida estándar (unitariamente invariante). Esta medida se usa para hallar volúmenes de conjuntos de estados.

\section{Idea conceptual}

Este nuevo tratamiento de los fundamentos de la mecánica estadística se tiene un universo. Dando como condición que el ambiente sea los suficientemente grande. Este universo está descrito por un estado cuántico puro (se conoce el estado de manera exacta) que obedece una restricción global. Se plantea que el sistema alcanza el equilibrio térmico por medio de la interacción mutua (termalización) es producto del entrelazamiento del sistema y el ambiente  \cite{Popescu2006}. Lo que se presentará más adelante es una definición más rigurosa que ayudará a dar cotas para la expresión "ambiente suficientemente grande".Esta idea permite formular un principio canónico general: el sistema estará termalizado para casi todos los estados puros del universo. esto es soportado por límites cuantitativos. La restricción que se impone no es una específica, esto generaliza los resultados tradicionales dados en la literatura donde se toma por restricción la energía \cite{KardarStat}.
\\
Ya poniendo lo dicho en un contexto un poco más matemático y siguiendo a Popescu et al. se supone tener un universo asilado y bastante grande este tiene dos partes el sistema $S$ y el ambiente $E$. La dimensión del ambiente es mucho más grande que la del sistema. Además se le impone una restricción global al universo llamada $R$. Desde la mecánica cuántica esto puede ponerse como restricciones en el espacio de Hilbert, restricción de los estados posibles:

\begin{equation}
\mathcal{H}_{R}\subseteq \mathcal{H}_{S}\otimes \mathcal{H}_{E},
\end{equation}

Donde $\mathcal{H}_{S}$ y $\mathcal{H}_{E}$ son los espacios de Hilbert del sistema y el ambiente con dimension $d_{S}$  y $d_{E}$ respectivamente. Es bueno recalcar que $R$ es una restricción arbitraria generalmente se toma como la energía del universo. Ahora se define el estado equiprobable del universo bajo $R$ como:

\begin{equation}
\mathcal{E}_{R} = \frac{1}{d_{R}} \mathbbm{1}_{R},
\end{equation}

Donde $\mathbbm{1}_{R}$ es el operador identidad (proyección) sobre el espacio de Hilbert  $\mathcal{H}_{R}$  que tiene dimensión $d_{R}$. Esto se relaciona con el principio de probabilidades iguales porque este es el estado máximamente mezclado en $\mathcal{H}_{R}$  \cite{SakuraiQuantum} por ser así todos los estados bajo la restricción de $R$ tienen la misma probabilidad de salir.
\\
Definimos $\Omega_{S}$ como el estado canónico que está restringido por $R$ cuando el universo se encuentra en el estado $\mathcal{E}_{R}$. Esto significa que si se hace una traza parcial del ambiente al universo da como resultado el estado canónico:

\begin{equation}
 \Omega_{S} =\Tr_{E} \mathcal{E}_{R}.
\end{equation}
\\
Para lo que sigue se hace un supuesto importante  y es que el universo está en un estado puro $\ket{\phi}$ y no en un estado mixto $\mathcal{E}_{R}$, esto quiere decir que se conoce todo lo que es permitido por la mecánica cuántica del universo. Si estuviese en un estado mixto significaría que nosotros no tenemos toda la información que se pudiese tener \cite{SakuraiQuantum}. Ahora lo que se quiere ver es que  pese a que el estado del universo es puro el estado reducido del sistema

\begin{equation}
\rho_{S}=\Tr_{E}\ket{\phi}\bra{\phi},
\end{equation}

se acerca al estado canónico para la gran mayoría de los casos es decir:

\begin{equation}
\rho_{S} \approx \Omega_{S}.
\end{equation}
\\
Por consiguiente para casi todos los estados puros del universo $\ket{\phi} \in \mathcal{H}_{R}$ el sistema se comporta como si el universo estuviese en el estado mixto equiprobable $\mathcal{E}_{R}$. Este  es el principio general canónico. Clarificando lo esbozado, el estado canónico del sistema $\Omega_{S}$ es el estado del sistema cuando el universo se encuentra en el estado equiprobable $\mathcal{E}_{R}$. Se puede interpretar el principio general canónico como un principio que estipula que las probabilidades iguales del sistema son aparentes porque para casi cualquier estado del universo, que sea puro, un subsistema de este universo que cumpla con ser lo suficientemente pequeño se comporta como si el universo estuviese en el estado equiprobable $\mathcal{E}_{R}$. Cabe recordar que aún no se ha especificado la restricción $R$ entonces todo este análisis es general, la restricción no necesariamente debe ser la energía u otras cantidades que se conserven. Esto hace que $\Omega_{S}$ no  deba ser obligatoriamente el estado canónico usual, puede ser el gran canónico o cualquier otro que sea acorde con la restricción impuesta\cite{ReichlStat}. Este principio puede ser de utilidad cuando la interacción entre el ambiente y el sistema no es débil o cuando las interacciones son complicadas el principio general canónico también aplica.

\section{Formulación matemática}

Hasta ahora no se han entrado en los detalles ni en qué significa bastante pequeño o bastante grandes, ni tampoco se ha demostrado el principio general canónico. En esta sección se orientará en los detalles matemáticos, las herramientas usadas y la demostración de los teoremas que Popescu et. al siguienron. En la siguiente sección se dará una perspectiva más física a lo hecho aquí.
\\
Para empezar se debe decir cuál será la distancia que usaremos para darle un sentido de cercanía a los estados $\rho_{S}$ y $\Omega_{S}$. La distancia a usar es una bastante conocida en la teoría cuántica de la información, la distancia de traza \cite{NielsenInformation}. Esta se define como:
\begin{equation}
D(\rho_{S}, \Omega_{S})= \frac{1}{2} \Tr |\rho_{S} -\Omega_{S}|=\frac{1}{2} \Tr \sqrt{(\rho_{S} -\Omega_{S})^{\dag}(\rho_{S} -\Omega_{S})}.
\end{equation}

Esta distancia de traza se relaciona con la distancia de norma de manera sencilla

\begin{equation}
\norm{\rho_{S}-\Omega_{S}}_{1}=2 D(\rho_{S}, \Omega_{S}).
\end{equation}

Teniendo ya una forma de darle sentido al concepto de que dos estados son cercanos entonces se puede plantear el teorema central de \cite{Popescu2006}:
\\
\begin{theorem} \label{teorema principal}
Para un estado escogido de manera aleatoria $\ket{\phi} \in \mathcal{H}_{R} \subseteq \mathcal{H}_{S} \otimes \mathcal{H}_{E} $ y un $\epsilon > 0$ arbitrario, la distancia entre la matriz densidad reducida del sistema $\rho_{S}=\Tr_{E}(\ket{\phi} \bra{\phi})$  y el estado canónico $\Omega_{S}=\Tr_{E} ( \mathcal{E}_{R})$ esta dado probabilísticamente por 

\begin{equation}
Pr_{\phi} \{  \norm{\rho_{S} -\Omega_{S}}_{1} \geq \eta \} \leq \eta',
\end{equation}

Donde 

\begin{equation}
\eta= \epsilon + \sqrt{ \frac{d_{S}}{d_{E}^{eff}} },
\end{equation}

\begin{equation}
\eta'=2\exp (-C d_{R} \epsilon^{2} ).
\end{equation}


y las constantes son: $ C=(18 \pi^{3})^{-1}, d_{R} = \dim \mathcal{H}_{R}, d_{S} = \dim \mathcal{H}_{S} $. $d_{E}^{eff}$ es la medida efectiva del tamaño del ambiente,
\begin{equation}
d_{E}^{eff}= \frac{1}{\Tr \Omega_{E}^{2}} \ge \frac{d_{R}}{d_{S}}.
\end{equation}

Donde $\Omega_{E}= \Tr_{S} \mathcal{E}_{R}$. Ambas cantidades $\eta $ y $\eta'$ serán pequeñas esto implica que el estado estará cercano al estado canónico con alta probabilidad cuando la dimensión efectiva del ambiente sea mucho más grande que la del sistema ( es decir $d_{E}^{eff} >> d_{S}$) y  $d_{R}\epsilon^2>>1>>\epsilon$. Esta última condición se puede asegurar cuando el espacio accesible total es grande ( es decir $d_{R}>>1$), escogiendo $\epsilon=d_{R}^{-\frac{1}{3}}$.\\
\end{theorem}

\subsection{Lema de Levy} \label{ levy}
Para poder demostrar el teorema \ref{teorema principal} se usará  el lema de Levy \cite{Lema}. Este dice que al seleccionar un punto $\phi$ aleatoriamente de una hiperesfera de dimensión alta y que $f(\phi)$ no cambie muy rápido, entonces $f(\phi) \approx \langle f \rangle $ con alta probabilidad, Más exactamente:\\

\begin{lemma} \label{lemma de levy}

Dada una función $f: \mathbb{S}^d \to \mathbb{R} $ definida en la hiperesfera d-dimensional $\mathbb{S}^d$ , y un punto $\phi \in \mathbb{S}^d $ es escogido de manera uniformemente aleatoria,
\begin{equation}
Pr_{\phi} \{ |f(\phi)- \langle f \rangle| \geq \epsilon \} \leq 2 \exp(-\frac{2C(d+1)\epsilon^2}{\eta^2})
\end{equation}
donde $\eta$ es la constante de Lipschitz de $f$, dado por $\eta= \sup|\nabla f|$ y $C=(18 \pi^3)^{-1} $.\\

\end{lemma}

Los conceptos manejados por el teorema \ref{lemma de levy} son conocidos generalmente a excepción de la llamada constante de Lipschitz.Para poder entender qué es la  constante de Lipschitz se debe ver primero qué significa que una función sea Lipschitz continua. 
\\
La definición de continuidad dad en el cálculo básico es:

\theoremstyle{definition}
\begin{definition}{continuidad}
Sea $f: I \to \mathbb{R}$ donde $I$ puede ser un intervalo abierto $(a,b)$ o uno cerrado $[a,b]$, además $C \in I$. Se dice que $f$ es continua en $C$ si y solo si para todo $ \epsilon >0 $  existe un $ \delta >0 $ tal que  $ |x-c|\longrightarrow |f(x)-f(c)|< \epsilon $.
\end{definition} 

La  definición anterior es la usada por lo general pero hay sutilezas en este concepto que no siempre son mostradas; como por ejemplo que $\delta$ depende de donde se ponga el  punto $C$, esto se ve claramente en la siguiente función: $f: (0,1) \to \mathbb{R}$, $f(x)=\frac{1}{x} $ al $C$ estar más lejos del $0$ permite un $\delta$ más grande pero al acercarse al $0$ el $\delta$ debe ser más pequeño. La continuidad de Lipschitz permite que se defina un $\delta$ constante sin importar donde se encuentre el $C$. Para resolver este detalle se motiva la definicón de Lipschitz continuo. Se es Lipschitz continuo con constante $\eta$ si

\begin{equation}
|f(x) -f(y)| \leq \eta |x-y|,
\end{equation}

esto permite decir que  $|f(x)-f(y)| < \epsilon $ entonces $\delta < \frac{\epsilon}{\eta} $. Ahora si f es derivable y $\nabla f$ es acotado, para $x$ y $y$ dados existe $\xi$ entre ellos tal que:


\begin{multline}
\implies f(x)-f(y)= \nabla f(\xi) (x-y) \\
\implies |f(x)- f(y)| \le |\nabla f(\xi)| |x-y| 
\implies |f(x)- f(y)| \le \sup|\nabla f(\xi)| |x-y|,
\end{multline}
como $\nabla f$ es acotado se tiene que $\sup |\nabla f(\xi)|=\eta$.\\

Gracias a la normalización, los estados puros en $\mathcal{H}_{R}$ se pueden representar como puntos sobre la superficie de una hiperesfera de dimensión $2d_{R}-1$, o sea  $\mathbb{S}^{2d_{R}-1}$ \cite{SakuraiQuantum}. Luego se puede aplicar \ref{lemma de levy} a estado cuánticos $\phi$ aleatoriamente seleccionados. Para los estados seleccionados aleatoriamente $\phi \in \mathcal{H}_{R}$, se desea mostrar que $\norm{\rho_{S}- \Omega_{S}}_{1} \approx 0$ con alta probabilidad. Para poder usar \ref{lemma de levy} primero  debe encontrarse la constante de Lipschitz ya teniendo una idea de qué es ser Lipschitz continuo se encontrará una cota para la constante $\eta$ de la función $f(\phi)=\norm{\rho_{S}-\Omega_{S}}_{1}$ que es la función que nos interesa para el problema físico. Entonces para poder lograr esto se procederá de la siguiente forma se definen dos estados reducidos $\rho_{1}= \Tr_{E} (\ket{\phi_{1}} \bra{\phi_{1}})$ y $\rho_{2}= \Tr_{E}(\ket{\phi_{2}} \bra{\phi_{2}})$, entonces

\begin{equation}
|f(\phi_{1})-f(\phi_{2})|^2= |\norm{\rho_{1}-\Omega}_{1} - \norm{\rho_{2}-\Omega}_{1}|^2.
\end{equation}

como $\norm{M}_{1}$ es una distancia (esto es $d(\rho_{1},\Omega)= \norm{\rho_{1}-\Omega}_{1}$) es cierto para un espacio métrico que

\begin{equation}
|d(x,z)-d(y,z)| \le d(x,y),
\end{equation}

Por lo tanto 

\begin{equation}
\Big | \norm{\rho_{1}-\Omega}_{1} - \norm{\rho_{2}-\Omega}_{1} \Big |^2 \le \norm{\rho_{1}- \rho_{2} }_{1}^2= \norm{\Tr_{E}(\ket{\phi_{1}} \bra{\phi_{1}}-\ket{\phi_{2}} \bra{\phi_{2}})}_{1}^{2}.
\end{equation}

Como existe una cota a la norma de una traza parcial dada por

\begin{equation}
\norm{\Tr_{\mathcal{B}}(M)}_{p} \le [dim(\mathcal{H_{B}})]^{\frac{p-1}{p}} \norm{M}_{p},
\end{equation}

entones la cota sobre los estados reducidos queda 
\begin{equation}
\norm{\Tr_{E}(\ket{\phi_{1}} \bra{\phi_{1}}-\ket{\phi_{2}} \bra{\phi_{2}})}_{1} \le \norm{\ket{\phi_{1}} \bra{\phi_{1}}-\ket{\phi_{2}} \bra{\phi_{2}}}_{1}.
\end{equation}

Por lo tanto se tiene hasta ahora:
\begin{equation}
\norm{\rho_{1}- \rho_{2} }_{1}^2 \le \norm{\ket{\phi_{1}} \bra{\phi_{1}}-\ket{\phi_{2}} \bra{\phi_{2}}}_{1}^{2},
\end{equation}

Usando la hermiticidad de $\rho$ y el teorema espectral se puede descomponer $\rho= UDU^{\dag}$ donde $U$ es un operador unitario y $D$ es diagonal. Junto con las propiedades de la traza $\Tr(\sqrt{UD^2U^{\dag}})= Tr(U\sqrt{D^2}U^{\dag})= \Tr(\sqrt{D^{2}})$ se llega a
\begin{equation}
\norm{\ket{\phi_{1}} \bra{\phi_{1}}-\ket{\phi_{2}} \bra{\phi_{2}}}_{1}^{2} = 4(1-|\bra{\phi_{1}} \ket{\phi_{2}}|^{2})
\end{equation}
entonces
\begin{equation}
 4(1-|\bra{\phi_{1}} \ket{\phi_{2}}|^{2}) \le 4|\ket{\phi}-\ket{\phi}|^2.
\end{equation}

Uniendo todos los pasos anteriores
\begin{equation}
|f(\phi_{1})-f(\phi_{2})|^{2} \le  4|\ket{\phi}-\ket{\phi}|^2
\end{equation}

o sea 
\begin{equation}
|f(\phi_{1})-f(\phi_{2})| \le  2|\ket{\phi}-\ket{\phi}|,
\end{equation}

Con esto se muestra que $\eta \le 2$.\\

\subsection{Demonstración del principio general canónico}
En esta parte se dará una demostración matemática explícita del principio general canónico usando el lema de Levy se entrará en los detalles matemáticos de usar este lema y las cotas adicionales que se necesitan para poder llegar al teorema \ref{teorema principal}. Habiendo dado una cota para $\eta$ ahora se puede usar por completo el lema \ref{lemma de levy} para la función $f(\phi)=\norm{\rho_{S} -\Omega_{S}}_{1}$ recordando que se reemplazará $d$ en el lema por $d=2d_{R}-1$.Tomando la parte derecha de la desigualdad de \ref{lemma de levy}  y sustituyendo la dimensión se tiene:

\begin{equation}
 2 \exp(-\frac{2C(d+1)\epsilon^2}{\eta^2})= 2 \exp(-\frac{4Cd_{R}\epsilon^2}{\eta^2})
\end{equation}

como $\eta \le 2$ entonces

\begin{equation} \label{desigualdad}
2\exp(-Cd_{R}\epsilon^2) \ge 2 \exp(-\frac{4Cd_{R}\epsilon^2}{\eta^2}) \ge \Pr_{\phi}[|f(\phi)- \langle f \rangle| \ge \epsilon ].
\end{equation}

Mirando más atentamente  $|f(\phi)- \langle f \rangle| \ge \epsilon $, como la norma de traza es un distancia se tiene que  $\norm{\rho_{S} -\Omega_{S}}_{1} \ge 0$ entonces


\begin{equation}
\norm{\rho_{S} -\Omega_{S}}_{1}- \langle \norm{\rho_{S} -\Omega_{S}}_{1} \rangle_{\phi} \ge \epsilon.
\end{equation}

Nombrando a $\mu = \epsilon + \langle \norm{\rho_{S} -\Omega_{S}}_{1} \rangle_{\phi} $ y $\mu'=2 \exp(-Cd_{R}\epsilon^2)$ esto permite organizar el lema de levy así:

\begin{equation} \label{eq.3}
\Pr_{\phi}[\norm{\rho_{S} -\Omega_{S}}_{1} \ge \mu] \le \mu'.
\end{equation}

Debido a que $d_{R}>>1$ Se asegura que $\epsilon$ y $\mu'$ son cantidades pequeñas al escoger $\epsilon=d_{R}^{-1/3}$. Para llegar a \ref{teorema principal} falta acotar $\mu$ con las dimensiones de los espacios conocidos, se impondrá una cota a $\langle \norm{\rho_{S} -\Omega_{S}}_{1} \rangle_{\phi}$; lo primero  para lograr esta empresa es acotar este promedio con trazas del estado del sistema  y luego se calcularán estas trazas para poder dejar la cota en términos de las dimensiones del sistema y la dimensión efectiva del ambiente. Se procede a encontrar la relación entre $\norm{\rho_{S} -\Omega_{S}}_{1}$ y $\norm{\rho_{S} -\Omega_{S}}_{2}$ esto se hará para tener una facilidad de manejo ya que la norma de Hilbert-Schmidt ($\norm{.}_{2}$) es más sencilla para trabajar que la norma de traza y luego se procederá con los planeado.
\\
La relación entre estas dos normas se puede ver desde el manejo de matrices. Sea $M$ una matriz $n \times n$ se sabe que si $M$ tiene $\lambda_{i}$ valores propios entonces:
\begin{equation}
\Tr M= \sum_{i} \lambda_{i}
\end{equation}
con esto se puede escribir de manera explícita la norma de traza 
\begin{equation}
\norm{M}_{1}^2=(\Tr|M|)^{2}=n^{2} \Big( \frac{1}{n}\sum_{i} |\lambda_{i}| \Big)^{2}.
\end{equation}
Como la función $x^{2}$ es convexa se puede usar la desigualdad de Jensen que dice: sean $a_{1},a_{2},...,a_{n} \le 0$ constantes y $a_{1} +...+a_{n}=1$ sea $f: I \to \mathbb{R}$ donde I es un intervalo, $x_{1},...,x_{n} \in I$ entonces
\begin{equation}
f(a_{1}x_{1}+...+a_{n}x_{n}) \le a_{1}f(x_{1})+...+a_{n}f(x_{n}),
\end{equation}

con esto se puede decir que

\begin{equation}
\Big( \frac{1}{n}\sum_{i} |\lambda_{i}| \Big)^{2} \le \frac{1}{n}\sum_{i} |\lambda_{i}|^{2}.
\end{equation}
Pero se sabe que 
\begin{equation}
\sum_{i} |\lambda_{i}|^2= \Tr(|M|^2)=\norm{M}_{2}^{2}
\end{equation}
se llega entonces a la conclusión que 
\begin{equation}
\norm{M}_{1}^{2} = n^{2}\Bigg( \frac{1}{n} \sum_{i}| \lambda_{i} | \Bigg)^{2} \le  n^{2} \frac{1}{n} \sum_{i} | \lambda_{i} |^{2}= n\norm{M}_{2}^{2}
\end{equation}
Gracias a lo anterior la relación entre normas es:
\begin{equation}\label{eq:0}
\norm{\rho_{S} - \Omega_{S}}_{1} \le \sqrt{d_{S}}\norm{\rho_{S} - \Omega_{S}}_{2}
\end{equation}
Esta relación se usará un poco más adelante.\\

Volviendo al cálculo de $\langle \norm{\rho_{S} -\Omega_{S}}_{1} \rangle_{\phi}$ se acotará la norma de Hilbert-Schmidt y con la relación entre normas se dará la desigualdad que limite el promedio de la norma de traza. Para empezar se recuerda que $ \langle f^{2} \rangle-\langle f \rangle^{2} \ge 0$ donde $\langle f \rangle = \int_{\mathcal{M}} f(x)p(x)dx$. Tomando a $f$ como $f= \norm{\rho_{S}-\Omega_{S}}_{2}$ entones

\begin{equation}
\langle \norm{ \rho_{S}-\Omega_{S} }_{2} \rangle \le \sqrt{ \langle \norm{ \rho_{S}-\Omega_{S} }_{2}^{2}\rangle}
\end{equation}

acordándose que este promedio es tomado a los estados $\ket{\phi}$, se omitirá por ahora el subíndice indicando este promedio, esto hace que $\Omega_{S}$ se tome constante. Por hermiticidad de $\rho_{S}-\Omega_{S}$ 

\begin{equation}
\sqrt{\langle \norm{ \rho_{S}-\Omega_{S} }_{2}^{2}\rangle} = \sqrt{\langle \Tr(\rho_{S}-\Omega_{S})^{2} \rangle}
\end{equation}

\begin{equation}
\sqrt{\langle \Tr( \rho_{S}-\Omega_{S})^{2} \rangle}= \sqrt{\langle \Tr(\rho_{S})^{2}\rangle -2 \Tr(\langle \rho_{S} \rangle \Omega_{S})+ \Tr(\Omega_{S}^{2}) }
\end{equation}

porque $\langle \rho_{S} \rangle=\Omega_{S}$. Luego se llega a

\begin{equation}
\langle \norm{ \rho_{S}-\Omega_{S} }_{2}\rangle \le \sqrt{\langle \norm{ \rho_{S}-\Omega_{S} }_{2}^{2}\rangle}= \sqrt{\langle \Tr(\rho_{S})^{2}\rangle - \Tr(\Omega_{S}^{2}) }.
\end{equation}

Por la relación entre la norma de traza y la norma de Hilbert-Schmidt se concluye lo que se quería

\begin{equation}
\langle \norm{ \rho_{S}-\Omega_{S} }_{1}\rangle \le \sqrt{d_{S}( \langle \Tr(\rho_{S})^{2}\rangle -\Tr(\Omega_{S}^{2}) ) }.
\end{equation}
\\
Aunque ya se ha acotado $\langle \norm{ \rho_{S}-\Omega_{S} }_{1}\rangle$ se quiere relacionar esta cota con las dimensiones del sistema para esto se procederá a demostrar la desigualdad
\begin{equation}
\langle \Tr \rho_{S}^{2} \rangle \le \Tr \langle \rho_{S} \rangle^{2} +\Tr \langle \rho_E \rangle^{2},
\end{equation}
recordando que el promedio es tomado con respecto a los estados $\ket{\phi}$, los métodos usados para encontrar esta desigualdad son usados también en destilación de entrelazamiento aleatoria y codificación de canal cuántico aleatorio \cite{QuantumDistilation}. Para poder hacer este cálculo se introduce una segunda copia del espacio de Hilbert. Ahora el problema se trabaja en $\mathcal{H}_{R} \otimes \mathcal{H}_{R'}$, donde $\mathcal{H}_{R'} \subseteq \mathcal{H}_{S'} \otimes \mathcal{H}_{E'}$.
Percatándose de lo siguiente

\begin{equation}
\Tr_{S} (\rho_{S})^2 = \sum_{k} (\rho_{kk})^{2} = \sum_{k,l,k',l'}(\rho_{kl})(\rho_{k'l'}) \bra{kk'}\ket{ll'}\bra{l'l'}\ket{kk'}.
\end{equation}

Sea $F_{SS'}$ la operación "flip" $S \longleftrightarrow S'$ definida de esta manera:

\begin{equation}
F_{SS'} =\sum_{S,S'} \ket{s'}\bra{s}_{S} \otimes \ket{s}\bra{s'}_{S'}
\end{equation}

Entonces

\begin{multline} \label{eq:1}
\sum_{k,l,k',l'}(\rho_{kl})(\rho_{k'l'}) \bra{kk'}\ket{ll'}\bra{l'l'}\ket{kk'}= \Tr_{SS'} ((\rho_{S} \otimes \rho_{S'})F_{SS'})\\
=\Tr_{RR'}((\ket{\phi} \bra{\phi} \otimes \ket{\phi} \bra{\phi})_{RR'}(F_{SS'} \otimes \mathbbm{1}_{EE'}))
\end{multline}

Pero como se quiere $\langle \Tr (\rho)^{2} \rangle= \int \Tr (\rho)^{2} d\phi $. Entonces para resolver esto se requiere saber $V=\int (\ket{\phi} \bra{\phi} \otimes \ket{\phi} \bra{\phi})d\phi$. V puede representarse como:

\begin{equation}
V= \alpha \Pi_{RR'}^{sim} + \beta \Pi_{RR'}^{anti},
\end{equation}
$\alpha$ y $\beta$ son constantes y $\Pi_{RR'}^{sim/anti}$ son proyectores en el subespacio simétrico y antisimétrico de $\mathcal{H_{R} \otimes H_{R'}}$ respectivamente, esto es posible por la invarianza unitaria de V. Debido a que 
\begin{equation}
(\ket{\phi} \bra{\phi} \otimes \ket{\phi} \bra{\phi})\frac{1}{\sqrt{2}} ( \ket{ab} -\ket{ba})=0  \quad  \forall a,b,\phi
\end{equation}
la parte antisimétrica siempre debe ser 0 entonces $\beta=0$. Por la normalización de $V$ se llega a $\alpha= \frac{1}{dim(RR'_{sim})}$, la dimensión del espacio $RR'_{sim}$ está dada por el álgebra lineal $dim(RR'_{sim})= \frac{d_{R}(d_{R}+1)}{2}$.Entonces
\begin{equation}
V= \langle \ket{\phi} \bra{\phi} \otimes \ket{\phi} \bra{\phi} \rangle = \frac{2}{d_{R}(d_{R} +1)} \Pi_{RR'}^{sim}.
\end{equation}
luego se puede escribir \ref{eq:1} como 
\begin{equation}
\langle \Tr (\rho)^{2} \rangle = \Tr_{RR'} \Big( \Big( \frac{2}{d_{R}(d_{R} +1)} \Pi_{RR'}^{sim} \Big)(F_{SS'} \otimes \mathbbm{1}_{EE'}) \Big)
\end{equation}
Al ser $\Pi_{RR'}^{sim}$ un proyector simétrico se escribe
\begin{equation}
\Pi_{RR'}^{sim}=\frac{1}{2}(\mathbbm{1}_{RR'}+(F_{RR'} ))
\end{equation}
donde $F_{RR'}$ es el operador "flip"  $R \longleftrightarrow R'$. Como $F_{RR'}$ es un operador que actúa sobre $RR'$ puede escribirse como $F_{RR'}= \mathbbm{1}_{RR'}(F_{SS'} \otimes F_{EE'})$.

reuniendo todo lo anterior

\begin{equation}
\langle \Tr_{S} \rho_{S}^2 \rangle = \Tr_{RR'} \Bigg( \frac{1}{d_{R}(d_{R}+1)} \Big( \mathbbm{1}_{RR'}+\mathbbm{1}_{RR'}(F_{SS'} \otimes F_{EE'}) \Big)  (F_{SS'}\otimes \mathbbm{1}_{EE'})   \Bigg)
\end{equation}

Distribuyendo y sabiendo que al hacer dos veces la operación "flip" no afecta nada se sigue lo siguiente

\begin{equation}
\Tr_{RR'} \Bigg( \frac{\mathbbm{1}_{RR'}}{d_{R}(d_{R}+1)} F_{SS'} \otimes \mathbbm{1}_{EE'} + \frac{\mathbbm{1}_{RR'}}{d_{R}(d_{R}+1)} (\mathbbm{1}_{SS'} \otimes F_{EE'})  \Bigg).
\end{equation}

Por las propiedades aditivas de la traza junto con $\mathbbm{1}_{R} \otimes \mathbbm{1}_{R'}$  y $\frac{1}{d_{R}(d_{R}+1)} \le \frac{1}{d_{R}^{2}}$ se llega a:

\begin{multline}
\langle \Tr_{S} \rho_{S}^2 \rangle= \Tr_{RR'} \Bigg( \Bigg( \frac{\mathbbm{1}_{RR'}}{d_{R}(d_{R}+1)} \Bigg) (F_{SS'} \otimes \mathbbm{1}_{EE'})  \Bigg) \\
+ \Tr_{RR'} \Bigg( \Bigg( \frac{\mathbbm{1}_{RR'}}{d_{R}(d_{R}+1)} \Bigg) (\mathbbm{1}_{SS'} \otimes F_{EE'} )  \Bigg)\\ 
\leq \Tr_{RR'} \Bigg( \Bigg( \frac{\mathbbm{1}_{R}}{d_{R}} \otimes  \frac{\mathbbm{1}_{R'}}{d_{R}}  \Bigg) (F_{SS'} \otimes \mathbbm{1}_{EE'})  \Bigg)\\
+ \Tr_{RR'} \Bigg( \Bigg( \frac{\mathbbm{1}_{R}}{d_{R}} \otimes  \frac{\mathbbm{1}_{R'}}{d_{R}}  \Bigg) (\mathbbm{1}_{SS'} \otimes F_{EE'} )  \Bigg)
\end{multline}

Recordando que $\frac{\mathbbm{1}_{R}}{d_{R}}= \mathcal{E}_{R}$ y $\Omega_{S}=\Tr_{E}(\mathcal{E}_{R})$ se tiene que:

\begin{multline}
\Tr_{RR'}\Bigg(\frac{\mathbbm{1}_{R}}{d_{R}} \otimes \frac{\mathbbm{1}_{R'}}{d_{R}} (F_{SS'} \otimes 		\mathbbm{1}_{EE'}) \Bigg) + \Tr_{RR'}\Bigg(\frac{\mathbbm{1}_{R}}{d_{R}} \otimes \frac{\mathbbm{1}_{R'}}{d_{R}} (\mathbbm{1}_{SS'}\otimes 	F_{EE'}) \Bigg) \notag \\
 = \Tr_{SS'}((\Omega_{S} \otimes \Omega_{S})F_{SS'}) + \Tr_{EE'}((\Omega_{E} \otimes \Omega_{E})F_{EE'})
\end{multline}
Por la ecuación \ref{eq:1} se llega a lo que se quería:	
\begin{equation}
\langle \Tr_{S} \rho_{S}^{2} \rangle \le \Tr_{S} \Omega_{S}^{2} +\Tr_{E} \Omega_{E}^{2},
\end{equation}
y esto es lo mismo que 
\begin{equation}
\langle \Tr_{S} \rho_{S}^{2} \rangle \le \Tr_{S} \langle \rho_{S} \rangle ^{2}  + \Tr_{E} \langle \rho_{E}.\rangle^{2}
\end{equation}

Gracias al resultado \ref{eq:0} se obtiene
\begin{equation}
\langle \norm{\rho_{S} -\Omega_{S}}_{1} \rangle \le \sqrt{d_{S}(Tr_{E} \langle \rho_{E} \rangle^{2})}.
\end{equation}
Si se define $d_{E}^{eff} \equiv \frac{1}{Tr_{E} \Omega_{E}^{2}}$ como la dimensión efectiva del ambiente en el estado canónico, esto mide la dimensión del espacio en el que el ambiente es más probable de estar, como $\langle \rho_{E} \rangle_{\phi} = \Omega_{E}$ se concluye que 
\begin{equation}
\langle \norm{\rho_{S} -\Omega_{S}}_{1} \rangle \le \sqrt{\frac{d_{S}}{d_{E}^{eff}}}
\end{equation}

Ya con esto se tiene el teorema \ref{teorema principal}. Cuando  el ambiente es mucho más grande que el sistema $\mu$ y $\mu'$  de la ecuación \ref{eq.3} serán pequeñas ($d_{E}^{eff}>>d_{S}$) implicando $\norm{\rho_{S} -\Omega_{S}}_{1} \approx 0$ con alta probabilidad. Aunque ya se llegó a la desigualdad que se quería se puede notar lo siguiente, sean los valores propios de $\Omega_{E}$ iguales a $\lambda_{E}^{k}$ con su máximo valor propio $\lambda_{E}^{max}$ se ve que 

\begin{multline}
\Tr_{E} \Omega_{E}^{2} = \sum_{k} (\lambda_{E}^{k})
\leq \lambda_{E}^{max} \sum_{k} \lambda_{E}^{k}\\
= \max_{\ket{\phi_{E}}} \bra{\phi_{E}} \Tr_{S} \Big( \frac{\mathbbm{1}_{R}}{d_{R}} \Big) \ket{\phi_{E}}\\
= \max_{\ket{\phi_{E}}} \sum_{s} \bra{s \phi_{E}} \frac{\mathbbm{1}_{R}}{d_{R}} \ket{s \phi_{E}} \leq \frac{d_{S}}{d_{E}}	
\end{multline}

en conclusión  $d^{eff}_{E} \geq d_{R}/ d_{S}$ entonces se obtiene

\begin{equation}
\langle \norm{\rho_{S} - \Omega_{S}}_{1} \rangle \leq \sqrt{\frac{d_{S}}{d_{E}^{eff}}} \leq  \sqrt{\frac{d_{S}^{2}}{d_{R}}}.
\end{equation}
Con esto se finaliza la demostración del teorema \ref{teorema principal} en la siguiente sección se hablará de sus consecuencias físicas.

\section{Significado físico}

El teorema anterior ya permite hablar del concepto importante que se sigue del aritculo de Popescu et al la idea general de la física es poder dar una relación uno a uno entre las propiedades de un objeto físico y su representación matemática con esta correspondencia se puede decir que la teoría esta completa cuando todas las propiedades que son posibles de medir tienen su semejante en la teoría \cite{Decoherence}. Por ejemplo en la física clásica a las cantidades como velocidad y distancia se le asigna los símbolos matemáticos $v$ y $x$, el cuerpo puede ser especificado dando estas dos cantidades en un tiempo determinado. La mecánica clásica nos da un ejemplo sencillo pero cuando se pasa a la mecánica estadística se encuentra perspectivas diferentes que entran en conflicto con esta sencilla idea. En los inicios de la mecánica estadística hubo muchas controversias dado a que ahora la idea determinista que Newton y otros habían mostrado como cierta (Se podía describir la naturaleza de manera predecible) ahora se introducía la probabilidad a sistemas que según Newton eran deterministas. En este punto se debe hacer claro que la necesidad de la probabilidad era por la falta de conocimiento pero no por un indeterminismo intrínseco como en la mecánica cuántica. Dos propuestas importantes surgieron en busca de una base conceptual sobre estas probabilidades, la visión de Gibbs y la de Boltzmann. Gibbs propuso la idea del ensamble, esta es una idea que hasta el día de hoy se sigue usando porque da los resultados correctos para las propiedades de sistemas termodinámicos.\\
La idea de Gibbs es la siguiente: Se tiene un sistema macroscópico con ciertos parámetros que se pueden medir. Este sistema está compuesto por varias partículas que tienen posiciones y momentos, pero estas propiedades están restringidas por los parámetros macroscópicos. Entonces como no se puede conocer el microestado del sistemas, no se puede medir cada posición y momento de cada partícula, se tiene una libertad sobre cuál es el microestado dado el macroestado. El ensamble es el conjunto de microestados que cumplen con el macroestado, mientras el microestado cumpla con los parámetros macroscópicos hará parte del ensamble así el sistema no se encuentre en ese microestado. En teoría no hay limitación sobre el número de microestados que pertenezcan a un ensamble, este puede tener infinitos microestados. Esta idea de base para la función de densidad en el espacio de fase, cada punto en este espacio es un microestado estos micro estados no interactuan entre ellos , con esto se puede construir la mecánica estadística que se conoce \cite{PathriaStat}. 
\\
Por el contrario Boltzmann da un concepción de la mecánica estadística más intuitiva; él propone $N$ partículas que interactúan entre ellas y cumplen los parámetros macroscópicos. Las partículas no son imaginarias como en el caso de los ensambles de Gibbs, ellas son las partículas que están en el sistema que se está estudiando. Se esperaría que la solución de Boltzmann fuera la que diera los resultados correctos dada la sencillez de su propuesta pero es la formulación de Gibbs la que da la termodinámica correcta. E.T Jaynes en \cite{JaynesEntropies} muestra las diferencias matemáticas de cada perspectiva y llega a la conclusión que la formulación de Gibbs da la entropía correcta sin importar cual sea el tipo de interacción que tenga el sistema, además encuentra que la entropía que sale de la formulación Boltzmanniana no es la misma entropía dada por la termodinámica,la entropía de Boltzmann es correcta cuando no hay interacción entre las partículas.
\\
Ya con la comprobación experimental se debería aceptar la perspectiva de Gibbs como la correcta y dejar a un lado la de boltzmann pero esto deja a la mecánica estadística en un contexto filosófico poco fuerte, aunque la física busque una correspondencia entre experimento y teoría no se puede aceptar cualquier teoría solo por tener congruencia con el experimento, la física debe buscar una estabilidad conceptual en sus teorías.¿ Por qué la perspectiva de Gibbs tiene un fundamento congruente ?  lo que propone Gibbs aunque matemáticamente aceptable se basa en todos los microestados en los que el macroestado no se encuentra es decir, que para poder determinar la propiedades de estado el cual se está estudiando se debe estudiar todos los estados en los que el sistema no se encuentra pero podría estar, pero por qué para analizar un sistema que se encuentra en un estado definido y exacto debería examinar los posibles infinitos estado en donde no se encuentra si se hace un promedio sobre estos microestados ¿ cuál es el significado de este promedio?. Este problema ontológico deja la perspectiva de Gibbs en duda, mientra la de Boltzmann habla del sistema que se examina sin ningún problema de este estilo. 
\\
La lucha entre las perspectiva ha sido una de aún no se resuelve lo que se ha mostrado hasta ahora ha sido la matemática de la tentativa de unir ambas perspectivas por Popescu et al. pero ¿ por qué estos resultado matemáticos intentan dar una luz a la resolución de estas perspectivas? se puede ver fácilmente que cuando se habla de del estado $ \Omega_{S} = \Tr_{B} \mathcal{E}_{R} $  este condensa la perspectiva de Gibbs mientras que $\rho_{S} = \Tr_{E} \ket{\phi} \bra{\phi}$ habla de un subsistema específico bajo una restricción esta es la perspectiva de Boltzmann por lo tanto el resultado mostrado acá dice: a menos que se tenga un estado bastante especial se puede asegurar que el subsistema de un universo muy grande, dado por un estado puro, se comporta de la misma manera que si el universo se encontrara en un estado máximamente mezclado. Esto corresponde a decir que para la mayoría de los estados en un espacio de Hilbert se puede reconciliar ambas perspectivas.
\\
Hay otro punto que da más mérito a la perspectiva mostrada acá y es que el estado $\Omega_{S}$ es la forma de mostrar la ignorancia subjetiva que se tiene del sistema. Se sabe que este es el punto de partida para la mecánica estadística pero esto también es un fundamento que generalmente es muy cuestionado gracias al teorema \ref{teorema principal} se cambia esa probabilidad subjetiva por una objetivad debido al entrelazamiento. Por este cambio a un fundamente más sólido no se debe entrar en el problema de ergodicidad ya que en general la solución del problema de esa probabilidad se dice que el sistema pasa por todos los estados,y eso solo complica la pregunta dado que hay sistemas que no pasan por todos los estados pero aún así son descritos por un estado canónico. Ya no debe entrarse en la pregunta sobre la ergodicidad del sistema dado que esta perspectiva muestra que los promedios no son necesarios en el momento de darle un fundamento a la mecánica estadística.





\thispagestyle{empty}
%\chapter{Termalización y los estados de Gibbs}
%\thispagestyle{empty}
\chapter{Evolución hacia el equilibrio}

Ya después de haber dado una explicación más sólida a los principios de la mecánica estadística la siguiente pregunta que debe responderse es,¿ cómo ocurre la termalización?. ¿ Esta termalización puede deducirse de las ecuaciones básicas ?. Este será la pregunta que Linden et al. quisieron atacar se procederá a mostrar sus resultados del artículo (Linden et al. ,2008). 
Una de tantas dificultades en la mecánica estadística es que sus principios se formulan desde la ignorancia subjetiva y los promedios de ensamble. Estos son principios físicos bastantes discutible.
Recientemente se ha dado cuenta que los promedio de ensamble y la ignorancia subjetiva no son necesarios, Porque sistemas cuánticos individuales pueden mostrar características estadísticas. Esto es debido al entrelazamiento lo cual es un efecto cuántico, esto hace que la falta de conocimiento ya no sea subjetiva sino objetiva. Por la teoría cuántica se sabe que aunque se tenga todo el sistema descrito por la función de onda, la cual da un conocimiento completo del sistema, los subsistemas de este pueden seguir en un estado desconocido. De forma más específica el sistema puede estar en un estado puro mientras que un subsistema puede estar en un estado mixto. En este caso no se puede conocer el subsistema ya que este se comporta como una densidad de probabilidad. La comparación con la contraparte clásica es sorprendente porque tener un conocimiento completo del sistema clásico es saber cualquier subsistema. Luego si se aplica una teoría clásica las probabilidades aparecen como falta del conocimiento posible de obtener, si se llega a conocer todo no abrían probabilidades.
Ya con resultados del capítulo anterior sobre: la mayoría de los estados puros de un sistema los subsistemas(suficientemente pequeños) están en un estado canónico. Este resultados aunque bastante general queda limitado por hablarse de un tiempo específico y habla de estados genéricos. Ahora se quiere ver su evolución temporal y en cuales circunstancias los sistemas llegan al equilibrio y fluctúan cerca a este. Los estados lejanos al equilibrio son no genéricos y serán tratados aquí. Para saber cómo moverse en este tema se va a determinar qué significa que un sistema esté termalizado. Para esto se darán las siguientes ideas:

\textbf{Equilibrio}: Se dice que un sistema se equilibra si este evoluciona a un estado específico(puede ser puro pero en general es mixto) y se mantiene allí por casi todo el tiempo. Dado esto no es importante cuál sea el estado de equilibrio este puede ser la distribución de Boltzmann o no. Además se puede relajar las condiciones de independencia del estado inicial, o sea puede depender del estado inicial del susbsistema y/o del estado inicial del ambiente de forma arbitraría. Esta definición de equilibrio es la parte más intuitiva a lo que se refiere sobre termalización. El hecho de que un estado se mantenga por un largo periodo de tiempo con las mismas características es lo que se piensa al  pensar en equilibrio. Aunque esto sería una forma general de hablar sobre equilibrio porque se da mucha libertad al estado y sus dependencias sobre el ambiente. Para seguir una idea más rigurosa de equilibrio se especifica lo sigeuinte.

\textbf{Independencia del ambiente}: El  estado de equilibrio del sistema no debería depender exactamente del estado inicial del baño. Esto quiere decir que el baño debe tener unos parámetros macroscópicos (como la temperatura) tales que al final llegan al equilibrio el estado dependa de la temperatura del baño. Con esto se puede restringir un poco más a lo que se llamará equilibrio.Esta idea también es proveniente de lo que normalmente se espera del equilibrio porque los parámetros macroscópicos son los que en general se tienen completamente especificados y se ve que para los mismos se tiene un mismo equilibrio. El estado exacto del baño no debería jugar un papel  tan importante ya que para los mismos parámetros macroscópicos pueden haber varios estados del baño que concuerden con ellos. Como se supone que con ciertos parámetros macroscópicos dados  un subsistema llegará al equilibrio sin importar cual ha sido su estado inicial. Pueden dos subsistemas preparados idénticamente haber sido iniciados en un estado A y el otro en el estado B, el equilibrio se obtendrá así hayan empezado de maneras diferentes. Esto motiva lo siguiente

\textbf{Independencia del estado del subsistema}:Si el subsistema es pequeño en comparación con el ambiente el estado de equilibrio del subsistema debería ser independiente de su estado inicial.

Pero para poder corroborar y unir estas ideas con los resultados ya bien conocidos se impone una última restricción

\textbf{Forma Boltzmanniana del estado de equilibrio}: Bajo condiciones del estado inicial y el hamiltoniano, el estado de equilibrio del subsistema puede ser escrito como $\rho_{S}= \frac{1}{Z} \exp(-\frac{H_{S}}{k_{B}T})$la forma familiar ya conocida.
Descomponiendo el problema de la termalización de esa forma permite ver cada uno de los aspectos por separado además de darle una generalidad a todo el tratamiento sin tener que restringirse a situaciones que usualmente se le asocian a la termalización. Por ejemplo no se debe quedar en el régimen de corta o débil interacción entre el sistema y el baño; decir que el baño es uno típico (Dado una temperatura o rango de energía), la energía puede llegar a ser una cantidad extensiva. Pueden tomarse situaciones en los que el sistema no llegue a equilibrio.
Se mostrará con suposiciones no muy fuertes que los primeros dos supuestos(Equilibrio e independencia del ambiente) son propiedades de sistemas cuánticos.



\textbf{El Hamiltoniano}: la evolución de todo el sistema viene dado por 
\begin{equation}
H= \sum_{k} E_{k} \ket{E_{k}} \bra{E_{k}},
\end{equation}
$\ket{E_{k}}$ es el estado propio con energía $E_{k}$.Al Hamiltoniano anterior se le dará la siguiente y única restricción: Que tenga brechas de energías no degeneradas. Esta restricción quiere decir que  un hamiltoniano tiene brechas de energía no degeneradas si para cualquier diferencia de energías propias que sea diferente acero determina los valores de energía involucrados. Por ejemplo si se tiene 4 valores propios de energía $E_{k}, E_{l},E_{m},E_{n}$ entonces $E_{k}-E_{l}=E_{m}-E_{n} $ implica $k=l$ y $m=n$, o $k=m$ y $l=n$. Esto implica que los niveles de energía no son iguales para diferentes estados (no son degenerados).
Esta restricción del Hamiltoniano implica que el subsistema y el baño no importa como se divida siempre van a estar interactuando. Esto excluye los hamiltonianos no interactuantes ($H=H_{S}+H_{E}$). Este tipo de hamiltonianos tienen muchas brechas de energía degeneradas. véase que si no hay interacción en el hamiltoniano la energía es $E=E_{S}+E_{E}$ sean $E_{1},E_{2},E_{3},E_{4}$ tales que se satisfaga $E_{i}=E_{i}^{S}+E_{i}^{E}$ $i=1,2,3,4$. Esto lleva a una brecha degenerada. Este supuesto no es tan fuerte como pueda llegar a parecer porque cualquier perturbación que se le haga al Hamiltoniano romperá las degeneraciones sin importar lo pequeña que seala perturbación. Aunque estos cambios se tarden en hacer efecto sobre la evolución del sistema las escalas temporales no son importantes en este momento. Con este supuesto se puede hablar de interacciones complejas que por lo general no son muy tratados como interacciones de larga distancia o interacciones entre todas las partículas esto hace que la energía no llegue a ser una cantidad extensiva.

\subsection{Equilibración}
El punto principal que se quiere dar a entender viene en forma del siguiente resultado: Para casa estado puro de un sistema cuántico que se compone por un número grande de estados de energía propios y el cual evoluciona bajo un Hamiltoniano que tiene brechas de energías no degeneradas y por lo demás arbitrario, es tal que cada subsistema pequeño llegará al equilibrio. Esto quiere decir que todos los subsistemas pequeños cumplirán con las ideas anteriores sobre equilibrio exactamente que el sistema evolucionará a un estado particular y se quedará cercano a él o en este durante la mayoría del tiempo.
Hay en este resultado un requerimiento que anteriormente no fue nombrado, el requerimiento de una cantidad grande de estados propios de energía. La necesidad de que el sistema tenga muchos estados propios de energía es equivalente a decir que el estado variará bastante durante su evolución temporal. viendo el caso trivial de un solo estado propio de energía es claro que este no cambiará para nada. Este caso tan particular no llega a ser de mucho interés porque en el sentido anterior de equilibrio no cambia para nada y se diría que este se encuentra ya en equilibrio. Los resultados que se quieren mostrar tomaŕan estados lejanos de equilibrio sistemas que no se encuentre en el pero aquí llegan las otras suposiciones hechas anteriormente para que el subsistema no dependa del estado inicial este debe perder la información entonces si el subsistema empieza lejano al equilibrio este pasará  por muchos estados en su camino al equilibrio lo cual implica que todo el sistema también evolucione en muchos estados. El hecho de que el subsistema haya llegado al equilibrio no significa que el sistema deje de evolucionar debido a la unitaridad este debe seguir evolucionando con la misma proporción de antes. Para que los estados en los que el subsistema se encuentre en no equilibrio ocurran poco, los estados del universo en los que el subsistema se encuentre en esas condiciones deben ser una fracción muy pequeña del total de estados por donde pasa el universo. Por esto el universo debe pasar por muchos estados y el requerimiento de que tenga muchos estados propios de energía se valida.
También se muestra que lo siguiente es suficiente: Cuando el estado de todo el sistema pasa por muchos estados diferentes cualquier estado pequeño del subsistema alcanza el equilibrio.
Otra forma de saber qué pasa con el sistema es observar qué ocurre con el ambiente. por unitaridad el sistema debe seguir evolucionando aunque el subsistema y se encuentre en equilibrio y no cambie. Esta Evolución puede darse por el cambio de correlaciones entre el subsistema y el baño o por cambios en el estado del baño. Lo que se muestra es: cuando el estado del baño pasa por muchos estados diferentes, cualquier subsistema alcanza el equilibrio.Con estas dos ideas se muestra que la equilibración ocurre en estados productos iniciales entre el subsistema y el baño, para casi todos los estados iniciales del baño.
La noción de evolución por muchos estados diferentes se puede escribir matemáticamente por la dimensión efectiva del estados promediado temporalmente $d^{eff}(\omega)$ donde $\omega=\langle \rho(t) \rangle_{t}$. La relación que se puede ver con los estados propios de energía es 
\begin{equation}
\ket{\psi(t)} = \sum_{k} c_{k} e^{-i\frac{E_{k}t}{\hbar}} \ket{E_{k}}
\end{equation}
luego el operador de densidad es 
\begin{equation}
\rho(t)=\sum_{k,l} c_{k} c_{l}^{*} e^{\frac{-i(E_{k}-E_{l})t}{\hbar}}\ket{E_{k}} \bra{E_{l}}
\end{equation}
entonces su promedio temporal es, recordando la condición de no degeneración de los niveles de energía
\begin{equation}
\omega= \sum_{k} |c_{k}|^{2} \ket{E_{k}} \bra{E_{k}}
\end{equation}
entoces la dimensión efectiva es 
\begin{equation}
d^{eff}(\omega)=\frac{1}{Tr(\omega^{2})}=\frac{1}{\sum_{k} |c_{k}|^{4} }
\end{equation}
De la misma manera el hechode que el baño pase por muchos estados diferentes viene dado por $d^{eff}(\omega_{B})$ por $\omega_{B}=\langle \rho(t) \rangle_{t}$. Como se espera que el baño al evolucionar pase por muchos más estados dado que el subsistema debe seguir evolucionando y el subsistema quede en un espacio de estados más pequeños se preve que $d^{eff}(\omega_{B})$ sea mucho más grande que $d_{S}$. Para formular ya el primer teorema se quiere ver la distancia entre $\rho_{S}(t)$ y su promedio temporal $\omega_{S}= \langle \rho_{S}(t) \rangle_{t}$. Como se espera que $\rho_{S}(t) $ vaya fluctuando alrededor de $\omega_{S}$ se analizará el promedio temporal de su distancia $\langle D(\rho_{S}(t) ,\omega_{S}) \rangle_{t}$ cuando este sea muy pequeño el subsistema debe pasar gran parte del tiempo muy cerca a $\omega_{S}$. Esto quiere decir que el subsistema se equilibrará(según la definición anterior) a $\omega_{S}$.

\textbf{Teorema 1}:Considere cualquier estado $\ket{\psi(t)} \in \mathcal{H}$ evolucionando bajo un Hamiltoniano con brechas de energía no-degeneradas. Luego la distancia promedio entre $\rho_{S}(t)$ y su promedio temporal $\omega_{S}$ está acotado por:

\begin{equation}
\langle D(\rho_{S}(t) ,\omega_{S}) \rangle_{t} \le \frac{1}{2} \sqrt{\frac{d_{S}}{d^{eff}(\omega-{B})}} \le \sqrt{\frac{d_{S}^{2}}{d^{eff}(\omega)}}.
\end{equation}
\textbf{Demostración}:Por la relación ya conocida entre la distanacia de traza y la norma de Hilbert-Schimidt que se usó en el capítulo anterior

\begin{equation}
\norm{M}_{1} \le \sqrt{n}\norm{M}_{2}
\end{equation}

se usa para la el operador $D(\rho_{1}, \rho_{2})$:

\begin{equation}
\frac{1}{2} Tr_{S} \sqrt{(\rho_{1} - \rho_{2})^{2}} \le \frac{1}{2}\sqrt{d_{S} Tr_{S} (\rho_{1}- \rho_{2})^{2}}
\end{equation}
Por la concavidad de la función raíz cuadrada

\begin{equation}
\langle D(\rho_{S}(t),\omega_{S}) \rangle_{t} \le \sqrt{  d_{S} \Big \langle Tr_{S}(\rho_{S}(t)-\omega_{S})^{2} \Big \rangle_{t}}
\end{equation}

usando las expansiones para $\rho_{S}$ y $\omega_{S}$
\begin{equation}
\rho_{S}(t)=\sum_{k,l} c_{k} c_{l}^{*} e^{\frac{-i(E_{k}-E_{l})t}{\hbar}} Tr_{B} (\ket{E_{k}} \bra{E_{l}}),
\end{equation}

\begin{equation}
\omega_{S} = \sum |c_{k}|^{2} Tr_{B}(\ket{E_{k}} \bra{E_{k}}),
\end{equation}

se puede escribir $\langle Tr_{S}(\rho_{S}(t)-\omega_{S})^{2} \rangle_{t}$ como 

\begin{equation}
\langle Tr_{S}(\rho_{S}(t)-\omega_{S})^{2} \rangle_{t}=\sum_{k \neq l} \sum_{m \neq n} \mathcal{T}_{klmn} Tr_{S} [Tr_{B}(\ket{E_{k}} \bra{E_{l}}) Tr_{B}(\ket{E_{k}} \bra{E_{l}})]
\end{equation}
donde $\mathcal{T}_{klmn}$ es :
\begin{equation}
\mathcal{T}_{klmn}=c_{k}c_{l}^{*}c_{m}c_{n}^{*} \Big \langle e^{\frac{-i(E_{k}-E_{l}+E_{m}-E_{n})t}{\hbar}} \Big \rangle_{t}
\end{equation}
debido a la restricción impuesta al Hamiltoniano de brechas de energías no degeneradas y como solo se toman elementos $k \neq l$ y $m\neq n$ los terminos que son diferentes de $0$ son $k=n$ y $l=m$.

Luego
 
\begin{multline}
\langle Tr_{S}(\rho_{S}(t)-\omega_{S})^{2} \rangle_{t} \\
	= \sum_{k \neq l} |c_{k}|^{2} |c_{l}|^{2} Tr_{S} [Tr_{B}(\ket{E_{k}} \bra{E_{l}}) Tr_{B}(\ket{E_{k}} \bra{E_{l}})] \\
	=\sum_{k \neq l} |c_{k}|^{2} |c_{l}|^{2} \sum_{ss' bb'} \bra{sb}\ket{E_{k}} \bra{E_{l}}\ket{s'b} \bra{s'b'}\ket{E_{l}} \bra{E_{k}}\ket{sb'} \\
	=\sum_{k \neq l} |c_{k}|^{2} |c_{l}|^{2} \sum_{ss' bb'} \bra{sb}\ket{E_{k}} \bra{E_{k}}\ket{sb'} \bra{s'b'}\ket{E_{l}} \bra{E_{l}}\ket{s'b}\\
	=\sum_{k \neq l} |c_{k}|^{2} |c_{l}|^{2} Tr_{B} [Tr_{S}(\ket{E_{k}} \bra{E_{k}}) Tr_{S}(\ket{E_{l}} \bra{E_{l}})] \\
	=\sum_{k \neq l}  Tr_{B} [Tr_{S}(|c_{k}|^{2}\ket{E_{k}} \bra{E_{k}}) Tr_{S}( |c_{l}|^{2} \ket{E_{l}} \bra{E_{l}})] 
\end{multline}
recordando la definición de $\omega_{B}$ y observando la segunda línea de la ecuación anterior
\begin{equation}
Tr_{B} \omega_{B}^{2} -\sum_{k} |c_{k}|^{2} Tr_{S}[ Tr_{B}( \ket{E_{k}} \bra{E_{k}})]^{2} \leq  Tr_{B}( \omega_{B}^{2}).
\end{equation}

Por la subaditividad débil de la entropía de Rényi:

\begin{equation}
Tr(\omega^{2}) \geq  \frac{Tr_{B}(\omega_{B}^{2})}{Rank(\rho_{S})} \geq \frac{Tr_{B}(\omega_{B}^{2})}{d_{S}}
\end{equation}

uniendo todos los resultados con la desigualdad inicial del promedio de la distancia

\begin{equation}
\langle D(\rho_{S}(t) ,\omega_{S}) \rangle_{t} \le  \frac{1}{2} \sqrt{d_{S} Tr_{B}(\omega_{B}^{2})} \le \frac{1}{2}\sqrt{d_{S}^{2}Tr(\omega^{2})}=\frac{1}{2}\sqrt{\frac{d_{S}^{2}}{d^{eff}(\omega)}}.
\end{equation}

Con este resultado se puede hablar de la termalización de una forma matemática consto se ve que el subsistema se equilibra cuando la dimensión de $d^{eff}(\omega)$ sea mucho mayor que la dimensión de dos copias del subsistema ($d_{S}^{2}$) o cuando la dimensión efectiva explorada por el baño $d^{eff}(\omega_{B})$ sea mucho más grande que la dimensión del subsistema.
EL resultado anterior tiene varias generalidades que se quiere recodar. La restricción impuesta sobre el Hamiltoniano  es una que no excluye muchos Hamiltonianos. Además de que esta ha sido la única restricción sobre todo el sistema , no se ha especificado nada del baño ni del subsistema. El baño no está necesariamente en equilibrio no se le ha dado ninguna interpretación térmica a ningún objeto tratado hasta ahora. no se ha hablado tampoco de ninguna forma en la que el subsistema llega al equilibrio ni que esta en algún estado específico.
Los valores propios de energía tampoco son importantes en las cotas dadas anteriormente, en el teorema al ser demostrado fue encontrado algunos valores propios de energía que al ser promediados dan 0. La energía es importante al buscar las formas exactas en las que evoluciona el sistema pero aquí se demostró que para la equilibración en intervalos de tiempo muy grandes no son muy importantes , las cotas son independientes del tiempo . La forma en que se dividió el sistema(subsistema y baño) solo es importante para el teorema 1 la dimensión del subsistema y no la especificación de la forma o un subsistema particular. Esto permite decir que cualquier subsistema con dimensión $d_{S}$ estará en equilibrio, los varios subsistemas bastante pequeños de dimensión $d_{S}$ también estarán en equilibrio.
\textbf{Equilibrio de conjuntos típicos}
Aunque se puso una cota a la fluctuación de $\rho_{S}(t)$ alrededor de $\omega_{S}$ se quiere ver cuales casos esta fluctuación es tan pequeña que el subsistema se equilibrará para esto se mostrará el siguiente teorema\\
\textbf{Terorema2}:
i)El promedio de la dimensión efectiva $\langle d^{eff}(\omega) \rangle_{\psi}$, donde el promedio es sobre estados puros aleatorios uniformemente distribuidos $\ket{\psi} \in \mathcal{H}_{R} \subset \mathcal{H}$. Es tal que 
\begin{equation}
\langle d^{eff}(\omega) \rangle_{\psi} \ge \frac{d_{R}}{2}
\end{equation}
ii)Para un estado aleatorio $\ket{\psi} \in \mathcal{H}_{R} \subset \mathcal{H}$, la probabilidad $Pr_{\psi} \{ d^{eff}(\omega) < \frac{d_{R}}{4}  \}$ de que $d^{eff}(\omega)$ es más pequeña que $\frac{d_{R}}{4}$ es:
\begin{equation}
Pr_{\psi} \{ d^{eff}(\omega) < \frac{d_{R}}{4}  \} \leq 2 \exp{-C \sqrt{d_{R}}}
\end{equation}
con constante $c= \frac{(ln 2)^{2}}{72 \pi^{3}} \approx 10^{-4}$.
La primera parte del teorema 2 nos habla de cómo el promedio de la dimensión efectiva es más grande que la dimensión del subespacio de Hilbert esto significa que si se tiene estados de un subespacio bastante grande se puede asegurar un $d^{eff}(\omega)$ grane que implica un $\langle D(\rho_{S}(t) ,\omega_{S}) \rangle_{t}$ pequeño. La segunda parte solo confirma de manera más estricta el hecho de encontrar un $d^{eff}(\omega)$ pequeño, mostrando que la probabilidad de encontrar una dimensión efectiva menor a $\frac{d_{R}}{4}$ es exponencialmente pequeña.\\
\textbf{equilibrio de estados genéricos}:
usando el teorema 2 para ver qué ocurre con un estado escogido de forma aleatorio del espacio total $\mathcal{H}$, un estado genérico. Esto es simplemente poner $d_{R}=d$ gracias a esto se comprende que $d^{eff}(\omega)~d$ ya que  hay una probabilidad exponencialmente baja para que se dé  el caso en que $d^{eff} < \frac{d}{4}$.Dado que $d=d_{S}d_{B}$ la cota para las fluctuaciones queda $\sqrt{\frac{d_{S}}{d_{B}}}$ para un sistema de muchas partículas la dimensión del espacio de Hilbert crece de manera exponencial, si el subsistema es una fracción constante del número de partículas del baño esta proporción caerá de manera exponencial con el número total de de partículas luego los subsistemas se equilibrarán.\\
\textbf{Equilibración de sistemas lejanos del equilibrio}:
Qué ocurre con los sistemas que están lejos del equilibrio. Lo dicho arriba ya cualquier se pensaría que  que cualquier sistema debe llegar al equilibrio pero esto no es cierto debido a que arriba se usó un estado genérico y los estados lejos del equilibrio no son genéricos; los estados lejos del equilibrio no son típicos por el capítulo anterior se sabe que cono cotas exponenciales la mayoría de los estados en el espacio de Hilbert son tales que un subsistema pequeño está en un estado canónico. Para poder sacar algo de esta pregunta se planteará la situación normal, hay un baño que consiste de un número muy grande de partículas de las cuales se conoce unos parámetros macroscópicos, dentro de este se pone un subsistema con un estado inicial arbitrario pero descorrelacionado con el ambiente. Ahora la pregunta es ¿ el subsistema se equilibra?, se verá que para cualquier estado inicial del subsistema y para casi todos los estados iniciales del baño el subsistema se equilibra. Esto incluye cuando el subsistema está lejos del equilibrio.
El estado inicial del sistema está dado por $\ket{\Psi}_{SB}=\ket{\psi}_{S} \ket{\psi}_{B}$. El estado del subsistema es uno arbitrario $\ket{\psi}_{S}$ en el espacio de Hilbert. Dado unos parámetros macroscópicos el baño  debe cumplir con estos luego el estado del baño $\ket{\phi}_{B} \in \mathcal{H}_{B}^{R} \subset\mathcal{H}_{B}$. Esta restricción mantiene la generalidad de seguir en cualquier espacio de Hilbert restringido. Puede tener o no un sentido termodinámico o macroscópico. Pero la restricción es solo inicial al evolucionar el baño en el tiempo este puede moverse fuera de $\mathcal{H}_{B}^{R}$. Usando el teorema 2 para $\mathcal{H}_{R}= \ket{\psi}_{S} \otimes \mathcal{H}_{B}^{R}$ luego $d_{R}=d_{B}^{R}$ esto da como resultado que para casi todos los estados iniciales del baño y cualquier estado del subsistema $d^{eff}(\omega) \geq \frac{d_{B}^{R}}{4}$ lo cual significa que el subsistema se equilibrará para esas condiciones, mientras que $d_{B}^{R} >> d_{S}^{2}$.

El mecanismo en que el subsistema se equilibra puede ser bastante complicada ya que el baño pasa por muchos estados diferentes y no llega el equilibrio. Aunque el baño no llegue a hacerlo y se salga del subespacio $\mathcal{H}_{B}^{R}$ y el subsistema  puede equilibrarse de todas formas. En principio puede que la evolución del subsistema sea sensible a la forma precisa del baño.
Para ver que el baño no se equilibra de manera genérica se verá que $d^{eff}(\omega_{S})$ es mucho mayor que $d^{eff}(\rho_{B}(t))$ lo cual muestra que el baño sigue evolucionando y no se equilibra en ningún estado. Como los dos sistemas están en un estado puro $Rank(\rho_{B}(t))=Rank(\rho_{S}(t)) \geq d_{S}$ como la dimensión efectiva de un estado es siempre menor a su rango se obtiene 

\begin{equation}
d^{eff}(\rho_{B}(t)) \geq d_{S}
\end{equation}

pero 

\begin{equation}
d^{eff}(\omega_{B})\leq \frac{d_{eff}(\omega)}{d_{S}}
\end{equation}

Pero para un estado genérico el teorema 2 ii dice $d^{eff}(\omega) > \frac{d_{R}}{4}$ se tiene

\begin{equation}
d^{eff}(\omega_{B}) \leq \frac{d^{R}}{4d_{S}}=\frac{d_{B}^{R}}{4d_{S}} >>d_{S} \leq d^{eff}(\rho_{B}(t))
\end{equation}
\textbf{Independencia del estado inicial}
LO que se ha mostrado hasta ahora ha sido como si el subsistema es pequeño este se equilibrará. Ahora el enfoque que se quiere mostrar es cual sería la dependencia del estado de equilibrio del subsistema. Hasta ahora el estado inicial podría hacer que el equilibrio sea un estado diferente dependiendo del inicial. Sea el estado de equilibrio del subsistema $\omega_{S}^{\psi}$.Como es conocido se desearía que el estado de equilibrio dependa solo de los parámetros macroscópicos y no del estado inicial microscópico.
El teorema siguiente prueba que para casi todos los estados en un subsistema restringido llevan al mismo estado de equilibrio.
\textbf{Teorema 3}:
i) Casi todos los estados iniciales de un subespacio restringido bastante grande llevan al mismo estado de equilibrio de un subsistema pequeño. En particula, con $\langle .\rangle_{\psi}$ siendo el promedio sobre estados puros aleatoriamente uniformes $\ket{\psi(0)} \in \mathcal{H}_{R} \subset \mathcal{H}_{S} \otimes \mathcal{H}_{B}$ u $\Omega_{S}= \langle \omega_{S}^{\psi} \rangle_{\psi}$ :
\begin{equation}
\langle D(\omega_{S}^{\psi}, \Omega_{S}) \rangle_{\psi} \leq \sqrt{\frac{d_{S} \delta}{4d_{R}}} \leq \sqrt{\frac{d_{S}}{4d_{R}}}
\end{equation}
La primera desigualdad es más estricta pero más complicada 
\begin{equation}
\delta= \sum_{k}\bra{E_{k}}\frac{\Pi_{R}}{d_{R}}\ket{E_{k}} Tr_{S} \Big( Tr_{B} (\ket{E_{k}} \bra{E_{k}}) \Big)^{2} \leq  1
\end{equation}
Donde $\Pi_{k}$ es el proyector sobre $\mathcal{H}_{R}$
ii)Paraunestado aleatorio $\ket{\psi} \in \mathcal{H}_{R}\subset \mathcal{H}$, la probabilidad que $D(\omega_{S}^{\psi},\Omega_{S}) > \frac{1}{2} \sqrt{\frac{d_{S} \delta}{d_{R}}}+ \epsilon$caiga exponencialmente con $\epsilon^{2}d_{R}$:
\begin{equation}
Pr_{\psi} \{ D(\omega_{S}^{\psi}, \Omega_{S}) > \frac{1}{2}\sqrt{\frac{d_{S} \delta}{d_{R}}} + \epsilon \} \leq 2exp(-C'\epsilon^{2}d_{R}),
\end{equation}
con $C'=\frac{2}{9 \pi^{3}}$ si se pone $\epsilon=d_{R}^{-1/3}$ da una distancia promedio pequeña con alta probabilidad cuando $d_{R}>>d_{S}$.

\textbf{Independencia del estado del ambiente}
PAra un estado inicial que sea el producto del estado del sistema y el del ambiente $\ket{\psi}_{SB}= \ket{\psi}_{S} \ket{\phi}_{B}$ en el espacio $\mathcal{H}_{R}=\ket{\psi} \otimes \mathcal{H}_{R}^{B}$.ya se mostró que para estados genéricos el ambiente causa que el subsistema se equilibre aunque el estado de equilibrio $\omega_{S}^{\psi}$ del subsistema podría  depender del estado inicial del baño $\ket{\phi}_{B}$. Esto no es así y casi todos los estados del baño en $\mathcal{H}_{B}^{R}$ llevan al mismo estado de equilibrio del subsistema simplemente usando el resultado anterior a  $\mathcal{H}_{B} = \ket{\psi}\otimes \mathcal{H}_{R}^{B}$ y luego $d_{R} = d_{R}^{B}$, si $d_{R}^{B}>>d_{S}$ para casi todos los estados iniciales del baño llevarán al subsistema al mismo estado de equilibrio $\omega_{S}$. como se usó la cota menos estricta este resultado no depende de la forma explícita de los estados de energías propios.
\textbf{Independenciadelestado del sistema}: Como yase hadicho el subsistema en contacto con el baño llega al equilibrio sin importar su estado inicial pero sí depende del baño. Este resultado no se pudo resolver pero se da más indicios.
Los problemas son que el estado de equilibrio no siempre es independiente del estado inicial del subsiste,a. si el subsistema puede cambiar al ambiente de forma drástica el estado de equilibrio sí dependerá del estado inicial del subsistema. aquí muestra que las dimensiones del espacio de Hilbert del subsistemas el ambiente son importante pero el valor se le haya dado una restricción al Hamiltoniano (Brechas de energías no degeneradas ) esta condición no están estricta como para no permitir cantidades conservadas en el subsistema cuando estas cantidades conservadas existen los estados del subsistema con diferentes constantes(estados inicial es diferentes )no podrán llegar al mismo estado de equilibrio.



\thispagestyle{empty}
\chapter{Conclusiones}

En este escrito se hizo un recuento histórico de la termodinámica para llegar a la mecánica estadística. Desde esta perspectiva histórica se mostró la evolución de los distintos conceptos que han fundamentado la mecánica estadística. Se expuso la perspectiva de Boltzmann junto con los las dificultades que esta acarrea. Como las objeciones dadas a su función H y la relación con la entropía termodinámica. La siguiente perspectiva que se exhibió fue la de Gibbs. La cual se condensa en la idea del ensamble. También se señaló los inconvenientes  que el concepto de copias imaginarias tiene, hablar de microestados que posiblemente pueden ser infinitos.\\
Junto con esto se muestra que hay otras perspectivas más modernas entre ellas la de Jaynes y su máxima entropía. La cual forma el primer vínculo entre la mecánica estadística y la teoría de la información. La perspectiva de Jaynes mejora respecto a las anteriores porque no introduce hipótesis a priori y especifica que la probabilidad es vista desde un sentido subjetivo, algo que las otras perspectivas no especificaban.
\\
Ya más adelante se siguió el artículo de Popescu et al  \cite{Popescu2006} mostrando la tipicidad canónica con ayuda de las herramientas de la teoría de la información cuántica. Se replicaron los cálculos del artículo para explicar los detalles del formalismo que los autores propusieron. Con esto puesto se vio la forma en que se conciliaban las perspectivas de Gibbs y Boltzmann gracias a este enfoque. Además el enfoque muestra que el principio de probabilidades a priori se puede reemplazar por el principio general canónico. En esto se concluye que es muy poco probable encontrar un estado que no se encuentre en el estado canónico.
\\
Para terminar se vio las extensiones que tiene los nuevos fundamentos gracias a la tipicidad canónica. Se vio la capacidad del principio general canónico tiene para poder explicar cómo se termaliza un sistema. Se reprodujo el análisis hecho por Linden et al \cite{LindenPaper} en este tema. Aunque no pudieron demostrar la independencia del subsistema, abrieron un campo a nuevas investigaciones.
\thispagestyle{empty}




%----------------------------------------------------------------------------------------
%	APÉNDICES
%----------------------------------------------------------------------------------------

\addtocontents{toc}{\vspace{2em}} % Agrega espacios en la toc

\appendix % Los siguientes capítulos son apéndices

%  Incluye los apéndices en el folder de apéndices

%\chapter{Álgebra Lineal}

Cosas.
%\thispagestyle{empty}
%\include{Apendices/AppendixB}
%\include{Apendices/AppendixC}

\addtocontents{toc}{\vspace{2em}} % Agrega espacio en la toc


%----------------------------------------------------------------------------------------
%	BIBLIOGRAFÍA
%----------------------------------------------------------------------------------------
%\backmatter
%\nocite{*}
%\bibliographystyle{unsrt}
%\bibliography{bibliografia} %Aquí ponen el nombre del archivo .bib

\begin{thebibliography}{10}
\bibitem{Popescu2006} Popescu, S., Short, A., \& Winter, A. (2006). Entanglement and the foundations of statistical mechanics. Nature Physics, 2(11), 754-758. http://dx.doi.org/ 10.1038/nphys444.
\bibitem{Goldstein} Goldstein, S., Lebowitz, J., Tumulka, R., \& Zanghì, N. (2006). Canonical Typicality. Physical Review Letters, 96(5). http://dx.doi.org/10.1103/physrevlett.96.050403
\bibitem{LandauStat} Landau, L., Lifshits, E., Pitaevskii, L., Sykes, J., \& Kearsley, M. (1980). Statistical physics. Amsterdam: Elsevier Butterworth Heinemann.
\bibitem{KardarStat}Kardar, M. (2007). cStatistical physics of particles. Cambridge: Cambridge University Press.
\bibitem{Pathria	Stat}Pathria, R., \& Beale, P. (2012). Statistical mechanics. Beijing: World Publishing Corporation.
\bibitem{HuangStat}Huang, K. (2005). Statistical mechanics. New York [u.a.]: Wiley.
\bibitem{Decoherence}Schlosshauer, M. (2010). Decoherence and the quantum-to-classical transition. Berlin: Springer.
\bibitem{CallenThermo}Callen, H. (1975). Thermodynamics. New York [etc.]: Wiley.
\bibitem{ReichlStat}Reichl, L. (2016). A Modern Course in Statistical Physics. Weinheim: Wiley.
\bibitem{SakuraiQuantum}Sakurai, J. (2013). Modern Quantum Mechanics. [S.l.]: Pearson Education India.
\bibitem{KhinchinStat}Khinchin, A., \& Gamow, G. Mathematical foundations of statistical mechanics.
\bibitem{TodaStat}Toda, M., Kubo, R., \& Saito, N. (1998). Statistical physics I. Berlin: Springer.
\bibitem{WildeInformation}Wilde, M. Quantum information theory.
\bibitem{NielsenInformation} Nielsen, M., \& Chuang, I. (2015). Quantum computation and quantum information. Cambridge: Cambridge University Press
\bibitem{LangAlgebra}Lang, S. (1972). Linear algebra. Reading, Mass.: Addison-Wesley.
\bibitem{VedralInformation}Vedral, V. (2013). Introduction to quantum information science. Oxford: Oxford University Press.
\bibitem{JaynesEntropies}Jaynes, E. (1965). Gibbs vs Boltzmann Entropies. American Journal Of Physics, 33(5), 391-398. http://dx.doi.org/10.1119/1.1971557
\bibitem{Lema}Vitali D. Milman., \& Gideon Schechtman. (1986). Asymptotic Theory of Finite Dimensional Normed Spaces. Springer Verlag.
\bibitem{QuantumDistilation}Horodecki, M., Oppenheim, J., \& Winter, A. (2005). Partial quantum information. Nature, 436(7051), 673-676. http://dx.doi.org/10.1038/nature03909.
\bibitem{LindenPaper}Linden, N., Popescu, S., Short, A., \& Winter, A. (2009). Quantum mechanical evolution towards thermal equilibrium. Physical Review E, 79(6). http://dx.doi.org/10.1103/physreve.79.061103.
\bibitem{RenyiEntropia}Hayden, P., van Dam, W.(2002). "Renyi-entropic bounds on quantum communication", arXiv:quant-ph/0204093.
\bibitem{MullerHistory} Muller, I. (2007). A history of thermodynamics. Berlin: Springer-Verlag GmbH.
\bibitem{Ehrenfest}Ehrenfest, P., Ehrenfest, T., \& Moravcsik, M. The conceptual foundations of the statistical approach in mechanics.
\bibitem{FermiPastaUlam} Fermi, E., Pasta, J., Ulam, S. (1955). Collected Papers of Enrico Fermi, Vol II, University of Chicago press, Chicago 1965.
\bibitem{Birkhoff}Birkhoff, G. (1931). Proof of the Ergodic Theorem. Proceedings Of The National Academy Of Sciences, 17(12), 656-660. http://dx.doi.org/10.1073/pnas.17.2.656
\bibitem{SusskindQuantum}Susskind, L., \& Friedman, A. (2015). Quantum mechanics. UK: Penguin Books.
\bibitem{ZettiliQuantum}Zettili, N. (2009). Quantum mechanics : concepts and applications.
\bibitem{JaynesI} Jaynes, E. (1957). Information Theory and Statistical Mechanics. Physical Review, 106(4), 620-630. http://dx.doi.org/10.1103/physrev.106.620
\bibitem{JaynesII} Jaynes, E. (1957). Information Theory and Statistical Mechanics. II. Physical Review, 108(2), 171-190. http://dx.doi.org/10.1103/physrev.108.171
\bibitem{ShannonInformation}Shannon, C. (1948). A Mathematical Theory of Communication. Bell System Technical Journal, 27(3), 379-423. http://dx.doi.org/10.1002/j.1538-7305.1948.tb01338.x
\bibitem{Penrose}Penrose, O. (2014). Foundations of Statistical Mechanics. Newburyport: Dover Publications.
\end{thebibliography}



 


\end{document}
